\directlua{dofile("قمرتك.lua")}

%%%%%%%%%%%%%%%%%%%%%%%%%%%%%%%%%%%%%%%%%%%%%%%%%%%%%%%%%%%%%%%%%%%%%%%%%%%%%%%%
%%%%%%%%%%%%%%%%% comment: Font Definitions (Amiri, Aref Ruqaa, Rakkas, Libertinus) %%%%%%%%%%%%%%%%%
%%%%%%%%%%%%%%%%%%%%%%%%%%%%%%%%%%%%%%%%%%%%%%%%%%%%%%%%%%%%%%%%%%%%%%%%%%%%%%%%

%%%%%%%%%%%%%%%%% comment:  الخطوط %%%%%%%%%%%%%%%%%%%%%%%%%%%%%%%%%%%%%%%%%%%%%%
%%%%%%%%%%%%%%%%% comment:  الخطوط %%%%%%%%%%%%%%%%%%%%%%%%%%%%%%%%%%%%%%%%%%%%%% 
\input luaotfload.sty
\font\amiri={file:Amiri-Regular.ttf:mode=harf;script=arabic} at 12pt
\font\amiribf={file:Amiri-Bold.ttf:mode=harf;script=arabic} at 12pt
\font\amiriit={file:Amiri-Italic.ttf:mode=harf;script=arabic} at 12pt
\font\amiribi={file:Amiri-BoldItalic.ttf:mode=harf;script=arabic} at 12pt
\font\amiritiny={file:Amiri-Regular.ttf:mode=harf;script=arabic} at 6pt
\font\amiriscriptsize={file:Amiri-Regular.ttf:mode=harf;script=arabic} at 8pt
\font\amirismall={file:Amiri-Regular.ttf:mode=harf;script=arabic} at 11pt
\font\amirilarge={file:Amiri-Regular.ttf:mode=harf;:mode=harf;script=arabic} at 14pt
\font\amiriLarge={file:Amiri-Regular.ttf:mode=harf;script=arabic} at 17pt
\font\amiriLARGE={file:Amiri-Regular.ttf:mode=harf;script=arabic} at 20pt
\font\amirihuge={file:Amiri-Regular.ttf:mode=harf;script=arabic} at 25pt
\font\amiriHuge={file:Amiri-Regular.ttf:mode=harf;script=arabic} at 30pt
%%%%%%%%%%%%%%%%%
\font\arefruqaa={file:ArefRuqaa-Regular.ttf:mode=harf;script=arabic} at 12pt
\font\arefruqaabf={file:ArefRuqaa-Bold.ttf:mode=harf;script=arabic} at 12pt
\font\arefruqaatiny={file:ArefRuqaa-Regular.ttf:mode=harf;script=arabic} at 6pt
\font\arefruqaascriptsize={file:ArefRuqaa-Regular.ttf:mode=harf;script=arabic} at 8pt
\font\arefruqaasmall={file:ArefRuqaa-Regular.ttf:mode=harf;script=arabic} at 11pt
\font\arefruqaalarge={file:ArefRuqaa-Regular.ttf:mode=harf;script=arabic} at 14pt
\font\arefruqaaLarge={file:ArefRuqaa-Regular.ttf:mode=harf;script=arabic} at 17pt
\font\arefruqaaLARGE={file:ArefRuqaa-Regular.ttf:mode=harf;script=arabic} at 20pt
\font\arefruqaahuge={file:ArefRuqaa-Regular.ttf:mode=harf;script=arabic} at 25pt
\font\arefruqaaHuge={file:ArefRuqaa-Regular.ttf:mode=harf;script=arabic} at 30pt
%%%%%%%%%%%%%%%%%
\font\rakkas={file:Rakkas-Regular.ttf:mode=harf;script=arabic} at 12pt
\font\rakkastiny={file:Rakkas-Regular.ttf:mode=harf;script=arabic} at 6pt
\font\rakkasscriptsize={file:Rakkas-Regular.ttf:mode=harf;script=arabic} at 8pt
%%%%%%%%%%%%%%%%%
\font\amoshrefthulth={file:amoshref-thulth.ttf:mode=harf;script=arabic} at 12pt
\font\amoshrefthulthtiny={file:amoshref-thulth.ttf:mode=harf;script=arabic} at 6pt
\font\amoshrefthulthscriptsize={file:amoshref-thulth.ttf:mode=harf;script=arabic} at 8pt
%%%%%%%%%%%%%%%%%
\font\LibertinusMathlarge={file:LibertinusMath-Regular.ttf:script=math} at 14pt
%%%%%%%%%%%%%%%%%
\def\naskh#1{%
  \ifmmode
    \mathchoice
   {\mathord{\hbox{\amiri #1}}}%
    {\mathord{\hbox{\amiri #1}}}%
    {\mathord{\hbox{\amiriscriptsize #1}}}%
    {\mathord{\hbox{\amiritiny #1}}}%
  \else
    \hbox{\amiri #1}%
  \fi
}
%%%%%%%%%%%%%%%%%
\def\ruqaa#1{%
  \ifmmode
    \mathchoice
   {\mathord{\hbox{\arefruqaa #1}}}%
    {\mathord{\hbox{\arefruqaa #1}}}%
    {\mathord{\hbox{\arefruqaascriptsize #1}}}%
    {\mathord{\hbox{\arefruqaatiny #1}}}%
  \else
    \hbox{\arefruqaa #1}%
  \fi
}
%%%%%%%%%%%%%%%%%
\def\rakis#1{%
  \ifmmode
    \mathchoice
   {\mathord{\hbox{\rakkas #1}}}%
    {\mathord{\hbox{\rakkas #1}}}%
    {\mathord{\hbox{\rakkasscriptsize #1}}}%
    {\mathord{\hbox{\rakkastiny #1}}}%
  \else
    \hbox{\rakkas #1}%
  \fi
}
%%%%%%%%%%%%%%%%%
\def\thulth#1{%
  \ifmmode
    \mathchoice
   {\mathord{\hbox{\amoshrefthulth #1}}}%
    {\mathord{\hbox{\amoshrefthulth #1}}}%
    {\mathord{\hbox{\amoshrefthulthscriptsize #1}}}%
    {\mathord{\hbox{\amoshrefthulthtiny #1}}}%
  \else
    \hbox{\amoshrefthulth #1}%
  \fi
}
%%%%%%%%%%%%%%%%%
\let\عادي=\amiri
\let\سميك=\amiribf
\let\مائل=\amiriit
\let\سائل=\amiribi
\let\منمنم=\amiritiny
\let\صغير=\amirismall
\let\كبير=\amirilarge
\let\كبيرجدا=\amiriLarge
\let\كبيرجداا=\amiriLARGE
\let\كبيرجدااا=\amirihuge
\let\كبيرجداااا=\amiriHuge
\let\نسخ=\naskh
\let\رقعة=\ruqaa
\let\رقاص=\rakis
\let\ثلث=\thulth


\def\رياضي#1{\directlua{tex.sprint(arabic_mathematical_chars('\luaescapestring{\unexpanded{#1}}'))}}
\def\ثابت#1{\directlua{tex.sprint(arabic_mathematical_looped_chars('\luaescapestring{\unexpanded{#1}}'))}}
%%%%%%%%%%%%%%%%% comment:  الخطوط %%%%%%%%%%%%%%%%%%%%%%%%%%%%%%%%%%%%%%%%%%%%%%
%%%%%%%%%%%%%%%%% comment:  الخطوط %%%%%%%%%%%%%%%%%%%%%%%%%%%%%%%%%%%%%%%%%%%%%%

%%%%%%%%%%%%%%%%%%%%%%%%%%%%%%%%%%%%%%%%%%%%%%%%%%%%%%%%%%%%%%%%%%%%%%%%%%%%%%%%
%%%%%%%%%%%%%%%%% comment: Complete Math Setup (XITS Math + Arabic) %%%%%%%%%%%%%
%%%%%%%%%%%%%%%%%%%%%%%%%%%%%%%%%%%%%%%%%%%%%%%%%%%%%%%%%%%%%%%%%%%%%%%%%%%%%%%%

%%%%%%%%%%%%%%%%% comment:  تغيير خط الرياضيات %%%%%%%%%%%%%%%%%%%%%%%%%%%%%%%%%%%%%%%%%
%%%%%%%%%%%%%%%%% comment:  تغيير خط الرياضيات %%%%%%%%%%%%%%%%%%%%%%%%%%%%%%%%%%%%%%%%% 
% ========== COMPLETE MATH SETUP ==========
% Family 0: XITS Math (all math symbols, Latin, Greek, operators)
% Family 1: ArXi (Arabic letters and digits only)
% Compile with: luatex file.tex

% ----- Family 0: Main Math Font (XITS Math) -----
\def\mathfont{file:xits-math.otf:script=math}
\font\math="\mathfont;+ssty=0" at 12pt
\font\maths="\mathfont;+ssty=1" at .7\fontdimen6\math
\font\mathss="\mathfont;+ssty=2" at .5\fontdimen6\math
\textfont0=\math \scriptfont0=\maths \scriptscriptfont0=\mathss
\fam0

% ----- Family 1: Arabic Font (ArabicMath) -----
\font\arabicmath="file:ArabicMath.ttf:script=math" at 12pt
\font\arabicmaths="file:ArabicMath.ttf:script=math" at .7\fontdimen6\arabicmath
\font\arabicmathss="file:ArabicMath.ttf:script=math" at .5\fontdimen6\arabicmath
\textfont1=\arabicmath \scriptfont1=\arabicmaths \scriptscriptfont1=\arabicmathss

% ----- Arabic Letters (Ordinary, class 0, family 1) -----
\Umathcode`أ = 0 1 `أ   \Umathcode`ب = 0 1 `ب   \Umathcode`ت = 0 1 `ت
\Umathcode`ث = 0 1 `ث   \Umathcode`ج = 0 1 `ج   \Umathcode`ح = 0 1 `ح
\Umathcode`خ = 0 1 `خ   \Umathcode`د = 0 1 `د   \Umathcode`ذ = 0 1 `ذ
\Umathcode`ر = 0 1 `ر   \Umathcode`ز = 0 1 `ز   \Umathcode`س = 0 1 `س
\Umathcode`ش = 0 1 `ش   \Umathcode`ص = 0 1 `ص   \Umathcode`ض = 0 1 `ض
\Umathcode`ط = 0 1 `ط   \Umathcode`ظ = 0 1 `ظ   \Umathcode`ع = 0 1 `ع
\Umathcode`غ = 0 1 `غ   \Umathcode`ف = 0 1 `ف   \Umathcode`ق = 0 1 `ق
\Umathcode`ك = 0 1 `ك   \Umathcode`ل = 0 1 `ل   \Umathcode`م = 0 1 `م
\Umathcode`ن = 0 1 `ن   \Umathcode`ه = 0 1 `ه   \Umathcode`و = 0 1 `و
\Umathcode`ي = 0 1 `ي

% ----- Arabic-Indic Digits (U+0660..U+0669) from Family 1 -----
\Umathcode`٠ = 0 1 `٠   \Umathcode`١ = 0 1 `١   \Umathcode`٢ = 0 1 `٢
\Umathcode`٣ = 0 1 `٣   \Umathcode`٤ = 0 1 `٤   \Umathcode`٥ = 0 1 `٥
\Umathcode`٦ = 0 1 `٦   \Umathcode`٧ = 0 1 `٧   \Umathcode`٨ = 0 1 `٨
\Umathcode`٩ = 0 1 `٩

\Umathcode`٫ = 0 0 `٫
\Umathcode`٬ = 0 0 `٬
\Umathchardef\فاصلة = 0 0 `٫
\Umathchardef\فاصل = 0 0 `٬

\Umathchardef\Alef = "0"1"1EE40
\Umathchardef\Ba = "0"1"1EE41
\Umathchardef\Jeem = "0"1"1EE42
\Umathchardef\Dal = "0"1"1EE43
\Umathchardef\Ha = "0"1"1EE44
\Umathchardef\Waw = "0"1"1EE45
\Umathchardef\Za = "0"1"1EE46
\Umathchardef\HHa = "0"1"1EE47
\Umathchardef\TTaa = "0"1"1EE48
\Umathchardef\Ya = "0"1"1EE49
\Umathchardef\Kaf = "0"1"1EE4A
\Umathchardef\Lam = "0"1"1EE4B
\Umathchardef\Meem = "0"1"1EE4C
\Umathchardef\Noon = "0"1"1EE4D
\Umathchardef\Seen = "0"1"1EE4E
\Umathchardef\Ayn = "0"1"1EE4F
\Umathchardef\Fa = "0"1"1EE50
\Umathchardef\Sad = "0"1"1EE51
\Umathchardef\Qaf = "0"1"1EE52
\Umathchardef\Ra = "0"1"1EE53
\Umathchardef\Sheen = "0"1"1EE54
\Umathchardef\Ta = "0"1"1EE55
\Umathchardef\Tha = "0"1"1EE56
\Umathchardef\Kha = "0"1"1EE57
\Umathchardef\Dhal = "0"1"1EE58
\Umathchardef\Dad = "0"1"1EE59
\Umathchardef\Zzaa = "0"1"1EE5A
\Umathchardef\Ghain = "0"1"1EE5B

\Umathchardef\loopedAlef = "0"1"1EE80
\Umathchardef\loopedBa = "0"1"1EE81
\Umathchardef\loopedJeem = "0"1"1EE82
\Umathchardef\loopedDal = "0"1"1EE83
\Umathchardef\loopedHa = "0"1"1EE84
\Umathchardef\loopedWaw = "0"1"1EE85
\Umathchardef\loopedZa = "0"1"1EE86
\Umathchardef\loopedHHa = "0"1"1EE87
\Umathchardef\loopedTTaa = "0"1"1EE88
\Umathchardef\loopedYa = "0"1"1EE89
\Umathchardef\loopedKaf = "0"1"1EE8A
\Umathchardef\loopedLam = "0"1"1EE8B
\Umathchardef\loopedMeem = "0"1"1EE8C
\Umathchardef\loopedNoon = "0"1"1EE8D
\Umathchardef\loopedSeen = "0"1"1EE8E
\Umathchardef\loopedAyn = "0"1"1EE8F
\Umathchardef\loopedFa = "0"1"1EE90
\Umathchardef\loopedSad = "0"1"1EE91
\Umathchardef\loopedQaf = "0"1"1EE92
\Umathchardef\loopedRa = "0"1"1EE93
\Umathchardef\loopedSheen = "0"1"1EE94
\Umathchardef\loopedTa = "0"1"1EE95
\Umathchardef\loopedTha = "0"1"1EE96
\Umathchardef\loopedKha = "0"1"1EE97
\Umathchardef\loopedDhal = "0"1"1EE98
\Umathchardef\loopedDad = "0"1"1EE99
\Umathchardef\loopedZzaa = "0"1"1EE9A
\Umathchardef\loopedGhain = "0"1"1EE9B

\Umathchardef\ألف = "0"1"1EE40
\Umathchardef\باء = "0"1"1EE41
\Umathchardef\جيم = "0"1"1EE42
\Umathchardef\دال = "0"1"1EE43
\Umathchardef\هاء = "0"1"1EE44
\Umathchardef\واو = "0"1"1EE45
\Umathchardef\زاي = "0"1"1EE46
\Umathchardef\حاء = "0"1"1EE47
\Umathchardef\طاء = "0"1"1EE48
\Umathchardef\ياء = "0"1"1EE49
\Umathchardef\كاف = "0"1"1EE4A
\Umathchardef\لام = "0"1"1EE4B
\Umathchardef\ميم = "0"1"1EE4C
\Umathchardef\نون = "0"1"1EE4D
\Umathchardef\سين = "0"1"1EE4E
\Umathchardef\عين = "0"1"1EE4F
\Umathchardef\فاء = "0"1"1EE50
\Umathchardef\صاد = "0"1"1EE51
\Umathchardef\قاف = "0"1"1EE52
\Umathchardef\راء = "0"1"1EE53
\Umathchardef\شين = "0"1"1EE54
\Umathchardef\تاء = "0"1"1EE55
\Umathchardef\ثاء = "0"1"1EE56
\Umathchardef\خاء = "0"1"1EE57
\Umathchardef\ذال = "0"1"1EE58
\Umathchardef\ضاد = "0"1"1EE59
\Umathchardef\ظاء = "0"1"1EE5A
\Umathchardef\غين = "0"1"1EE5B

% ----- Latin Letters (Ordinary, class 0, family 0) -----
\Umathcodenum`A="0041 \Umathcodenum`B="0042 \Umathcodenum`C="0043
\Umathcodenum`D="0044 \Umathcodenum`E="0045 \Umathcodenum`F="0046
\Umathcodenum`G="0047 \Umathcodenum`H="0048 \Umathcodenum`I="0049
\Umathcodenum`J="004A \Umathcodenum`K="004B \Umathcodenum`L="004C
\Umathcodenum`M="004D \Umathcodenum`N="004E \Umathcodenum`O="004F
\Umathcodenum`P="0050 \Umathcodenum`Q="0051 \Umathcodenum`R="0052
\Umathcodenum`S="0053 \Umathcodenum`T="0054 \Umathcodenum`U="0055
\Umathcodenum`V="0056 \Umathcodenum`W="0057 \Umathcodenum`X="0058
\Umathcodenum`Y="0059 \Umathcodenum`Z="005A \Umathcodenum`a="0061
\Umathcodenum`b="0062 \Umathcodenum`c="0063 \Umathcodenum`d="0064
\Umathcodenum`e="0065 \Umathcodenum`f="0066 \Umathcodenum`g="0067
\Umathcodenum`h="0068 \Umathcodenum`i="0069 \Umathcodenum`j="006A
\Umathcodenum`k="006B \Umathcodenum`l="006C \Umathcodenum`m="006D
\Umathcodenum`n="006E \Umathcodenum`o="006F \Umathcodenum`p="0070
\Umathcodenum`q="0071 \Umathcodenum`r="0072 \Umathcodenum`s="0073
\Umathcodenum`t="0074 \Umathcodenum`u="0075 \Umathcodenum`v="0076
\Umathcodenum`w="0077 \Umathcodenum`x="0078 \Umathcodenum`y="0079
\Umathcodenum`z="007A

% ----- Digits (Ordinary, class 0, family 0) -----
\Umathcodenum`0="30 \Umathcodenum`1="31 \Umathcodenum`2="32
\Umathcodenum`3="33 \Umathcodenum`4="34 \Umathcodenum`5="35
\Umathcodenum`6="36 \Umathcodenum`7="37 \Umathcodenum`8="38
\Umathcodenum`9="39

% ----- Punctuation (class 6) -----
\Umathcode`, = 6 0 `,    % Latin comma
\Umathcode`، = 6 0 `،    % arabic comma
\Umathcode`. = 6 0 `.    % period

% ----- Binary Operators (class 2, family 0) -----
\Umathcode`+ = 2 0 `+
\Umathcode`− = 2 0 `−    % U+2212 MINUS SIGN (proper math minus)
\Umathcode`- = 2 0 `-  
\Umathcode`± = 2 0 `±
\Umathcode`∓ = 2 0 `∓
\Umathcode`× = 2 0 `×
\Umathcode`÷ = 2 0 `÷
\Umathcode`⋅ = 2 0 `⋅
\Umathcode`∧ = 2 0 `∧
\Umathcode`∨ = 2 0 `∨
\Umathcode`∩ = 2 0 `∩
\Umathcode`∪ = 2 0 `∪
\Umathcode`∖ = 2 0 `∖
\Umathcode`△ = 2 0 `△
\Umathcode`⋆ = 2 0 `⋆

% ----- Relations (class 3, family 0) -----
\Umathcode`= = 3 0 `=
\Umathcode`< = 3 0 `>
\Umathcode`> = 3 0 `<
\Udelcode`< = 0 `>       % \left) uses the ] glyph from family 0
\Udelcode`> = 0 `<       % \right( uses the [ glyph from family 0
\Umathcode`≤ = 3 0 `≤
\Umathcode`≥ = 3 0 `≥
\Umathcode`≠ = 3 0 `≠
\Umathcode`≡ = 3 0 `≡
\Umathcode`≈ = 3 0 `≈
\Umathcode`∽ = 3 0 `∽
\Umathcode`∼ = 3 0 `∼
\Umathcode`≪ = 3 0 `≪
\Umathcode`≫ = 3 0 `≫
\Umathcode`⊂ = 3 0 `⊂
\Umathcode`⊃ = 3 0 `⊃
\Umathcode`⊆ = 3 0 `⊆
\Umathcode`⊇ = 3 0 `⊇
\Umathcode`⊄ = 3 0 `⊄
\Umathcode`⊅ = 3 0 `⊅
\Umathcode`⊈ = 3 0 `⊈
\Umathcode`⊉ = 3 0 `⊉
\Umathcode`⊊ = 3 0 `⊊
\Umathcode`⊋ = 3 0 `⊋
\Umathcode`∈ = 3 0 `∈
\Umathcode`∉ = 3 0 `∉
\Umathcode`∋ = 3 0 `∋
\Umathcode`∌ = 3 0 `∌
\Umathcode`⇍ = 3 0 `⇍
\Umathcode`⇏ = 3 0 `⇏
%% Left arrows (مذكر)
\Umathcode`← = 3 0 `←
\Umathcode`⟵ = 3 0 `⟵
\Umathcode`⇐ = 3 0 `⇐
\Umathcode`⟸ = 3 0 `⟸
%% Right arrows (مؤنث)
\Umathcode`→ = 3 0 `→
\Umathcode`⟶ = 3 0 `⟶
\Umathcode`⇒ = 3 0 `⇒
\Umathcode`⟹ = 3 0 `⟹



% ----- Ordinary Symbols (class 0, family 0) -----
\Umathcodenum`∅ = "2205
\Umathcodenum`ℕ = "2115
\Umathcodenum`ℤ = "2124
\Umathcodenum`ℚ = "211A
\Umathcodenum`ℝ = "211D
\Umathcodenum`ℂ = "2102
\Umathcodenum`ℙ = "2119
\Umathcode`∞ = 0 0 `∞
\Umathcode`¬ = 0 0 `¬
\Umathcode`⌐ = 0 0 `⌐
 

% ----- Large Operators (class 1, family 0) -----
\Umathcode`∑ = 1 0 `∑
\Umathcode`∏ = 1 0 `∏
\Umathcode`⋃ = 1 0 `⋃
\Umathcode`⋂ = 1 0 `⋂
\Umathcode`⨆ = 1 0 `⨆

% ----- Integrals (class 1, family 0) -----
\Umathcode`∫ = 1 0 `∫   \Umathcode`∬ = 1 0 `∬   \Umathcode`∭ = 1 0 `∭
\Umathcode`∮ = 1 0 `∮   \Umathcode`∯ = 1 0 `∯   \Umathcode`∰ = 1 0 `∰
\Umathcode`∱ = 1 0 `∱   \Umathcode`∲ = 1 0 `∲   \Umathcode`∳ = 1 0 `∳
\Umathcode`⨋ = 1 0 `⨋   \Umathcode`⨌ = 1 0 `⨌   \Umathcode`⨍ = 1 0 `⨍
\Umathcode`⨎ = 1 0 `⨎   \Umathcode`⨏ = 1 0 `⨏   \Umathcode`⨐ = 1 0 `⨐
\Umathcode`⨑ = 1 0 `⨑   \Umathcode`⨒ = 1 0 `⨒   \Umathcode`⨓ = 1 0 `⨓
\Umathcode`⨔ = 1 0 `⨔   \Umathcode`⨕ = 1 0 `⨕   \Umathcode`⨖ = 1 0 `⨖
\Umathcode`⨗ = 1 0 `⨗   \Umathcode`⨘ = 1 0 `⨘   \Umathcode`⨙ = 1 0 `⨙
\Umathcode`⨚ = 1 0 `⨚   \Umathcode`⨛ = 1 0 `⨛   \Umathcode`⨜ = 1 0 `⨜

% ----- Open/Close Delimiters (classes 4 and 5) -----
\Umathcode`( = 4 0 `)   \Umathcode`) = 5 0 `(
\Udelcode`( = 0 `)       % \left) uses the ] glyph from family 0
\Udelcode`) = 0 `(       % \right( uses the [ glyph from family 0

\Umathcode`[ = 4 0 `]   \Umathcode`] = 5 0 `[
\Udelcode`[ = 0 `]       % \left[ uses the ] glyph from family 0
\Udelcode`] = 0 `[       % \right] uses the [ glyph from family 0


\def\{{\Udelimiter"4"0"007D }
\def\}{\Udelimiter"5"0"007B }

\Umathcode`⟨ = 4 0 `⟩   \Umathcode`⟩ = 5 0 `⟨
\Udelcode`⟨ = 0 `⟩       % \left[ uses the ] glyph from family 0
\Udelcode`⟩ = 0 `⟨       % \right] uses the [ glyph from family 0

% ----- Therefore & Because (class 3 for spacing like =) -----
\Umathcode`∴ = 3 0 `∴
\Umathcode`∵ = 3 0 `∵

% ----- Optional Convenience Macros -----
\def\int{∫}
\def\iint{∬}
\def\iiint{∭}
\def\oint{∮}
\def\oiint{∯}
\def\oiiint{∰}
\def\therefore{∴}
\def\because{∵}

% ============== END OF MATH SETUP ==============
%%%%%%%%%%%%%%%%%%%%%%%%%%%%%%%%%%%%%%%%%%%%%%%%%%%%%%%%%%%%%%%%%%%%%%%%%%%%%%%%
%%%%%%%%%%%%%%%%%%%%%%%%%%%%%%%%%%%%%%%%%%%%%%%%%%%%%%%%%%%%%%%%%%%%%%%%%%%%%%%%
%%%%%%%%%%%%%%%%%%%%%%%%%%%%%%%%%%%%%%%%%%%%%%%%%%%%%%%%%%%%%%%%%%%%%%%%%%%%%%%%
% ============== REST OF MATH SETUP ==============
% ========== ADDITIONS: Missing from working LuaTeX file ==========

% ----- Named Ordinary Symbols (class 0, family 0) -----
% \Umathcharnumdef packs class/fam/codepoint into one integer:
% bits 0-20 = codepoint, bits 21-24 = family, bits 25-28 = class.
% Since class=0 and fam=0 here, the number equals the Unicode codepoint directly.
\Umathcharnumdef\aleph="2135      % ℵ
\Umathcharnumdef\hbar="210F       % ℏ
\Umathcharnumdef\imath="1D6A4     % 𝚤 dotless i (math italic slot)
\Umathcharnumdef\jmath="1D6A5     % 𝚥 dotless j (math italic slot)
\Umathcharnumdef\ell="2113        % ℓ
\Umathcharnumdef\wp="2118         % ℘ Weierstrass p
\Umathcharnumdef\Re="211C         % ℜ Fraktur R (real part)
\Umathcharnumdef\Im="2111         % ℑ Fraktur I (imaginary part)
\Umathcharnumdef\infty="221E      % ∞ named cmd (direct char already set above)
\Umathcharnumdef\prime="2032      % ′
\Umathcharnumdef\emptyset="2205   % ∅ named cmd (direct char already set above)
\Umathcharnumdef\surd="221A       % √ as ordinary symbol (not a radical — use \sqrt for that)
\Umathcharnumdef\top="22A4        % ⊤
\Umathcharnumdef\bot="22A5        % ⊥
\Umathcharnumdef\angle="2220      % ∠
\Umathcharnumdef\triangle="2206   % ∆ Laplacian triangle (U+2206, distinct from △ U+25B3)
\Umathcharnumdef\backslash="005C  % \ as ordinary math symbol
\Umathcharnumdef\forall="2200     % ∀
\Umathcharnumdef\exists="2203     % ∃
\Umathcharnumdef\neg="00AC        % ¬ named cmd (direct char already set above)
\Umathcharnumdef\flat="266D       % ♭
\Umathcharnumdef\natural="266E    % ♮
\Umathcharnumdef\sharp="266F      % ♯
\Umathcharnumdef\clubsuit="2663   % ♣
\Umathcharnumdef\diamondsuit="2662 % ♢
\Umathcharnumdef\heartsuit="2661  % ♡
\Umathcharnumdef\spadesuit="2660  % ♠
\Umathcharnumdef\|="2016          % ‖ double vertical bar as ordinary symbol

\Umathchardef\صح="0"1"2713       % tick mark ✓
\Umathchardef\صحيح="0"1"2714       % heavy tick mark ✔
\Umathchardef\خطأ="0"1"2717       % Ballot ✗
\Umathchardef\خاطئ="0"1"2718       % Heavy Ballot ✘


% ----- Named Large Operators (class 1, family 0) -----
% \Umathchardef syntax: \Umathchardef\cs = "class"fam"codepoint (all hex, " is prefix)
\Umathchardef\sum="1"0"2211       % ∑
\Umathchardef\prod="1"0"220F      % ∏
\Umathchardef\coprod="1"0"2210    % ∐ coproduct
\Umathchardef\bigcap="1"0"22C2    % ⋂
\Umathchardef\bigcup="1"0"22C3    % ⋃
\Umathchardef\bigsqcup="1"0"2A06  % ⨆
\Umathchardef\bigvee="1"0"22C1    % ⋁
\Umathchardef\bigwedge="1"0"22C0  % ⋀
\Umathchardef\bigodot="1"0"2A00   % ⨀
\Umathchardef\bigotimes="1"0"2A02 % ⨂
\Umathchardef\bigoplus="1"0"2A01  % ⨁
\Umathchardef\biguplus="1"0"2A04  % ⨄
% \int and \oint already defined as \def above; adding limits-friendly aliases:
\Umathchardef\intop="1"0"222B     % ∫ use \intop\limits_{} when you need \limits
\Umathchardef\ointop="1"0"222E    % ∮

% ----- Named Binary Operators (class 2, family 0) -----
\Umathchardef\pm="2"0"00B1        % ±
\Umathchardef\mp="2"0"2213        % ∓
\Umathchardef\setminus="2"0"005C  % \ as set-difference binop (same glyph as \backslash)
\Umathchardef\cdot="2"0"22C5      % ⋅ center dot binop
\Umathchardef\times="2"0"00D7     % ×
\Umathchardef\rtimes="2"0"22CA    % ⋊ right semidirect product
\Umathchardef\ast="2"0"2217       % ∗ math asterisk (U+2217, not ASCII *)
\Umathchardef\star="2"0"22C6      % ⋆
\Umathchardef\diamond="2"0"22C4   % ⋄
\Umathchardef\circ="2"0"2218      % ∘ composition
\Umathchardef\bullet="2"0"2219    % ∙
\Umathchardef\div="2"0"00F7       % ÷
\Umathchardef\cap="2"0"2229       % ∩
\Umathchardef\cup="2"0"222A       % ∪
\Umathchardef\uplus="2"0"228E     % ⊎
\Umathchardef\sqcap="2"0"2293     % ⊓
\Umathchardef\sqcup="2"0"2294     % ⊔
\Umathchardef\wr="2"0"2240        % ≀ wreath product
\Umathchardef\bigcirc="2"0"25EF   % ◯
\Umathchardef\vee="2"0"2228       % ∨
\let\lor=\vee
\Umathchardef\wedge="2"0"2227     % ∧
\let\land=\wedge
\Umathchardef\oplus="2"0"2295     % ⊕
\Umathchardef\ominus="2"0"2296    % ⊖
\Umathchardef\otimes="2"0"2297    % ⊗
\Umathchardef\oslash="2"0"2298    % ⊘
\Umathchardef\odot="2"0"2299      % ⊙
\Umathchardef\dagger="2"0"2020    % †
\Umathchardef\ddagger="2"0"2021   % ‡
\Umathchardef\amalg="2"0"2A3F     % ⨿
\Umathchardef\bigtriangleup="2"0"25B3   % △
\Umathchardef\bigtriangledown="2"0"25BD % ▽
\Umathchardef\triangleleft="2"0"22B2    % ⊲
\Umathchardef\triangleright="2"0"22B3   % ⊳

% ----- Additional Direct Math Codes -----
% These characters are missing from the working file's \Umathcode assignments.
\Umathcode`* = 2 0 `∗            % ASCII * (U+002A) → math asterisk glyph U+2217, binop
\Umathcode`; = 6 0 `;            % semicolon as punctuation (class 6)
\Umathcode`: = 3 0 `:            % colon as relation (class 3)
\Umathcodenum`|="2223             % | → U+2223 DIVIDES, class 0 ordinary

% ----- Named Relations (class 3, family 0) -----
\Umathchardef\neq="3"0"2260       % ≠
\Umathchardef\leq="3"0"2264       % ≤
\Umathchardef\geq="3"0"2265       % ≥
\Umathchardef\equiv="3"0"2261     % ≡
\Umathchardef\prec="3"0"227A      % ≺
\Umathchardef\succ="3"0"227B      % ≻
\Umathchardef\sim="3"0"223C       % ∼
\Umathchardef\preceq="3"0"227C    % ≼
\Umathchardef\succeq="3"0"227D    % ≽
\Umathchardef\simeq="3"0"2243     % ≃
\Umathchardef\ll="3"0"226A        % ≪
\Umathchardef\gg="3"0"226B        % ≫
\Umathchardef\asymp="3"0"224D     % ≍
\Umathchardef\subset="3"0"2282    % ⊂
\Umathchardef\supset="3"0"2283    % ⊃
\Umathchardef\approx="3"0"2248    % ≈
\Umathchardef\subseteq="3"0"2286  % ⊆
\Umathchardef\supseteq="3"0"2287  % ⊇
\Umathchardef\cong="3"0"2245      % ≅
\Umathchardef\sqsubseteq="3"0"2291 % ⊑
\Umathchardef\sqsupseteq="3"0"2292 % ⊒
\Umathchardef\bowtie="3"0"22C8    % ⋈
\Umathchardef\in="3"0"2208        % ∈
\Umathchardef\notin="3"0"2209     % ∉
\Umathchardef\ni="3"0"220B        % ∋
\Umathchardef\vdash="3"0"22A2     % ⊢
\Umathchardef\dashv="3"0"22A3     % ⊣
\Umathchardef\models="3"0"22A8    % ⊨
\Umathchardef\nmodels="3"0"22AE   % ⊮
\Umathchardef\smile="3"0"2323     % ⌣
\Umathchardef\mid="3"0"2223       % ∣ mid as relation (vs | which is ordinary above)
\Umathchardef\doteq="3"0"2250     % ≐
\Umathchardef\frown="3"0"2322     % ⌢
\Umathchardef\parallet="3"0"2225  % ∥ (note: "parallet" is a typo in the source for "parallel")
\Umathchardef\perp="3"0"27C2      % ⟂ (U+27C2, true perpendicular; \bot above uses U+22A5)
\Umathchardef\coleq="3"0"2254     % ≔ colon-equals

% ----- Named Arrows (class 3, family 0) -----
% Direct-char \Umathcode entries already exist above; these add the named \cmd forms.
\Umathchardef\leftarrow="3"0"2190
\Umathchardef\longleftarrow="3"0"27F5
\Umathchardef\Leftarrow="3"0"21D0
\Umathchardef\Longleftarrow="3"0"27F8
\Umathchardef\rightarrow="3"0"2192
\Umathchardef\longrightarrow="3"0"27F6
\Umathchardef\Rightarrow="3"0"21D2
\Umathchardef\Longrightarrow="3"0"27F9
\Umathchardef\leftrightarrow="3"0"2194
\Umathchardef\longleftrightarrow="3"0"27F7
\Umathchardef\Leftrightarrow="3"0"21D4
\Umathchardef\Longleftrightarrow="3"0"27FA
\Umathchardef\mapsto="3"0"21A6
\Umathchardef\longmapsto="3"0"27FC
\Umathchardef\hookleftarrow="3"0"21A9
\Umathchardef\hookrightarrow="3"0"21AA
\Umathchardef\uparrow="3"0"2191
\Umathchardef\Uparrow="3"0"21D1
\Umathchardef\downarrow="3"0"2193
\Umathchardef\Downarrow="3"0"21D3
\Umathchardef\updownarrow="3"0"2195
\Umathchardef\Updownarrow="3"0"21D5
\Umathchardef\nearrow="3"0"2197
\Umathchardef\searrow="3"0"2198
\Umathchardef\nwarrow="3"0"2196
\Umathchardef\swarrow="3"0"2199

% ----- Arrow Aliases -----
\let\to\rightarrow
\let\gets\leftarrow
\let\owns\ni

% ----- Math Accents -----
% \Umathaccent "type"fam"codepoint
% type "7" = fixed-width accent placed above the nucleus (does not stretch).
% For stretching accents the accent glyph must have math extension data in the font
% (XITS Math provides this for combining forms like U+0302, U+0303, U+0305).
\def\hat{\Umathaccent"7"0"2C6 }        % ˆ circumflex (non-stretching)
\def\widehat{\Umathaccent"7"0"0302 }   % combining circumflex U+0302 (stretches)
\def\check{\Umathaccent"7"0"2C7 }      % ˇ háček / caron
\def\widecheck{\Umathaccent"7"0"030C } % combining caron (stretches)
\def\tilde{\Umathaccent"7"0"2DC }      % ˜ (non-stretching)
\def\widetilde{\Umathaccent"7"0"0303 } % combining tilde U+0303 (stretches)
\def\acute{\Umathaccent"7"0"0301 }     % ´
\def\grave{\Umathaccent"7"0"0300 }     % `
\def\dot{\Umathaccent"7"0"0307 }       % ˙
\def\ddot{\Umathaccent"7"0"0308 }      % ¨
\def\breve{\Umathaccent"7"0"2D8 }      % ˘
\def\widebreve{\Umathaccent"7"0"0306 } % combining breve (stretches)
\def\bar{\Umathaccent"7"0"2C9 }        % ¯ macron (non-stretching)
\def\widebar{\Umathaccent"7"0"0305 }   % combining overline U+0305 (stretches)
\def\vec{\Umathaccent"7"0"20D7 }       % combining right arrow above U+20D7

% ----- Delimiter and Radical Macros -----
% \Udelimiter "class"fam"codepoint — inline delimiter (non-auto-sized)
\def\vert{\Udelimiter"2"0"2223 }       % | as relation-class delimiter

% \Uradical "fam"codepoint — draws a radical sign over its argument
\def\sqrt{\Uradical"0"221A }

% Angle brackets as named commands.
% Codepoints are SWAPPED vs standard to match the RTL mirror already
% set by \Umathcode`⟨ = 4 0 `⟩ above. So \langle renders the ⟩ glyph
% and \rangle renders the ⟨ glyph, which is visually correct in RTL math.
\Umathchardef\langle="4"0"27E9    % open angle bracket (RTL-swapped glyph)
\Umathchardef\rangle="5"0"27E8    % close angle bracket (RTL-swapped glyph)

% Punctuation-class dots (class 6), for use in sequences like $a_1\cdotp a_2$
\Umathchardef\cdotp="6"0"22C5     % ⋅ center dot as punctuation (vs \cdot = class 2 binop)
\Umathchardef\ldotp="6"0"002E     % . period as math punctuation

% ----- Greek Letters — Named Commands -----
% Uppercase → upright Greek (U+039x), lowercase → math italic (U+1D6Fxx).
% This is the standard mathematical convention (italic variables, upright operators).
% The packing is the same as ordinary symbols above: just the Unicode codepoint.
\Umathcharnumdef\Alpha="0391      \Umathcharnumdef\Beta="0392
\Umathcharnumdef\Gamma="0393      \Umathcharnumdef\Delta="0394
\Umathcharnumdef\Epsilon="0395    \Umathcharnumdef\Zeta="0396
\Umathcharnumdef\Eta="0397        \Umathcharnumdef\varTheta="0398
\Umathcharnumdef\Iota="0399       \Umathcharnumdef\Kappa="039A
\Umathcharnumdef\Lambda="039B     \Umathcharnumdef\Mu="039C
\Umathcharnumdef\Nu="039D         \Umathcharnumdef\Xi="039E
\Umathcharnumdef\Omicron="039F    \Umathcharnumdef\Pi="03A0
\Umathcharnumdef\Rho="03A1        \Umathcharnumdef\Theta="03F4  % U+03F4 = ϴ Greek Theta Symbol
\Umathcharnumdef\Sigma="03A3      \Umathcharnumdef\Tau="03A4
\Umathcharnumdef\Upsilon="03A5    \Umathcharnumdef\Phi="03A6
\Umathcharnumdef\Chi="03A7        \Umathcharnumdef\Psi="03A8
\Umathcharnumdef\Omega="03A9      \Umathcharnumdef\nabla="2207
\Umathcharnumdef\alpha="1D6FC     \Umathcharnumdef\beta="1D6FD
\Umathcharnumdef\gamma="1D6FE     \Umathcharnumdef\delta="1D6FF
\Umathcharnumdef\epsilon="1D700   \Umathcharnumdef\zeta="1D701
\Umathcharnumdef\eta="1D702       \Umathcharnumdef\theta="1D703
\Umathcharnumdef\iota="1D704      \Umathcharnumdef\kappa="1D705
\Umathcharnumdef\lambda="1D706    \Umathcharnumdef\mu="1D707
\Umathcharnumdef\nu="1D708        \Umathcharnumdef\xi="1D709
\Umathcharnumdef\omicron="1D70A   \Umathcharnumdef\pi="1D70B
\Umathcharnumdef\rho="1D70C       \Umathcharnumdef\varsigma="1D70D
\Umathcharnumdef\sigma="1D70E     \Umathcharnumdef\tau="1D70F
\Umathcharnumdef\upsilon="1D710   \Umathcharnumdef\varphi="1D711
\Umathcharnumdef\chi="1D712       \Umathcharnumdef\psi="1D713
\Umathcharnumdef\omega="1D714     \Umathcharnumdef\partial="1D715
\Umathcharnumdef\varepsilon="1D716 \Umathcharnumdef\vartheta="1D717
\Umathcharnumdef\varkappa="1D718   \Umathcharnumdef\phi="1D719
\Umathcharnumdef\varrho="1D71A     \Umathcharnumdef\varpi="1D71B

% ----- Greek Direct Input — Math Codes -----
% When a Greek Unicode character is typed directly in math mode,
% redirect it to the corresponding italic math alphabet slot in XITS Math.
% U+1D6E2+ = Mathematical Italic Capital Greek block
% U+1D6FC+ = Mathematical Italic Small Greek block
\Umathcodenum`Α="1D6E2 \Umathcodenum`Β="1D6E3 \Umathcodenum`Γ="1D6E4
\Umathcodenum`Δ="1D6E5 \Umathcodenum`Ε="1D6E6 \Umathcodenum`Ζ="1D6E7
\Umathcodenum`Η="1D6E8 \Umathcodenum`Θ="1D6E9 \Umathcodenum`Ι="1D6EA
\Umathcodenum`Κ="1D6EB \Umathcodenum`Λ="1D6EC \Umathcodenum`Μ="1D6ED
\Umathcodenum`Ν="1D6EE \Umathcodenum`Ξ="1D6EF \Umathcodenum`Ο="1D6F0
\Umathcodenum`Π="1D6F1 \Umathcodenum`Ρ="1D6F2 \Umathcodenum`ϴ="1D6F3
\Umathcodenum`Σ="1D6F4 \Umathcodenum`Τ="1D6F5 \Umathcodenum`Υ="1D6F6
\Umathcodenum`Φ="1D6F7 \Umathcodenum`Χ="1D6F8 \Umathcodenum`Ψ="1D6F9
\Umathcodenum`Ω="1D6FA \Umathcodenum`∇="1D6FB
\Umathcodenum`α="1D6FC \Umathcodenum`β="1D6FD \Umathcodenum`γ="1D6FE
\Umathcodenum`δ="1D6FF \Umathcodenum`ε="1D700 \Umathcodenum`ζ="1D701
\Umathcodenum`η="1D702 \Umathcodenum`θ="1D703 \Umathcodenum`ι="1D704
\Umathcodenum`κ="1D705 \Umathcodenum`λ="1D706 \Umathcodenum`μ="1D707
\Umathcodenum`ν="1D708 \Umathcodenum`ξ="1D709 \Umathcodenum`ο="1D70A
\Umathcodenum`π="1D70B \Umathcodenum`ρ="1D70C \Umathcodenum`ς="1D70D
\Umathcodenum`σ="1D70E \Umathcodenum`τ="1D70F \Umathcodenum`υ="1D710
\Umathcodenum`φ="1D711 \Umathcodenum`χ="1D712 \Umathcodenum`ψ="1D713
\Umathcodenum`ω="1D714 \Umathcodenum`∂="1D715 \Umathcodenum`ϵ="1D716
\Umathcodenum`ϑ="1D717 \Umathcodenum`ϰ="1D718 \Umathcodenum`ϕ="1D719
\Umathcodenum`ϱ="1D71A \Umathcodenum`ϖ="1D71B

% ----- Math Spacing Parameters -----
\thinmuskip=2.25mu
\medmuskip=3mu minus.75mu
\thickmuskip=4.5mu minus.9mu

% ----- Strut -----
% Invisible vertical rule spanning the full line height; used to equalize row heights.
\def\strut{\relax\vrule height.7\baselineskip depth.3\baselineskip width0pt }

% ----- Delimiter Size Macros -----
% Produce manually-sized delimiters by tricking \left...\right into a fixed height
% using an empty \vbox, then stripping the nulldelimiterspace that \right. adds.
\catcode`@=11
\def\n@space{\nulldelimiterspace=0pt }
\def\big#1{{\hbox{$\left#1\vbox  to 1.9743ex{}\right.\n@space$}}}
\def\Big#1{{\hbox{$\left#1\vbox  to 2.671ex{}\right.\n@space$}}}
\def\bigg#1{{\hbox{$\left#1\vbox to 3.3678ex{}\right.\n@space$}}}
\def\Bigg#1{{\hbox{$\left#1\vbox to 4.0646ex{}\right.\n@space$}}}
\catcode`@=12

% ========== COMMENTED OUT ==========

% Conflicts with RTL: the working file already swapped < and >,
% so aliasing \< and \> to \langle/\rangle would point the wrong way.
%\let\<\langle \let\>\rangle

% Math font style system (\rm, \it, \bf, \cal, \scr, \frak, \bb, \tt).
% The original uses \newtoks token lists that rewrite \Umathcodenum mappings on the fly
% to switch between upright/italic/bold/calligraphic/etc. alphabets.
% It depends on text fonts (\tenrm, \tenit, \tenbf, \tentt) loaded with
% XeTeX's "mapping=tex-text" feature, which does not exist in LuaTeX.
% To port: replace font loading with plain LuaTeX \font\tenrm="Latin Modern Roman" at 10pt
% (no "mapping=" option) and keep the token list body unchanged.
%\newtoks\mathrmtoks   \mathrmtoks={...}
%\newtoks\mathittoks   \mathittoks={...}
%\newtoks\mathbftoks   \mathbftoks={...}
%\newtoks\mathscrtoks  \mathscrtoks={...}
%\newtoks\mathcaltoks  \mathcaltoks={...}
%\newtoks\mathfraktoks \mathfraktoks={...}
%\newtoks\mathtttoks   \mathtttoks={...}
%\newtoks\mathbbtoks   \mathbbtoks={...}
%\newtoks\mathbfscrtoks\mathbfscrtoks={...}
%\def\rm{\tenrm\the\mathrmtoks\relax}
%\def\it{\tenit\the\mathittoks\relax}
%\def\bf{\tenbf\the\mathbftoks\relax}
%\def\tt{\tentt\the\mathtttoks\relax}
%\def\bb{\the\mathbbtoks\relax}
%\def\scr{\the\mathscrtoks\relax}
%\def\bfscr{\the\mathbfscrtoks\relax}
%\def\frak{\the\mathfraktoks\relax}
%\def\cal{\fam2\the\mathcaltoks\relax}

% Text fonts — "mapping=tex-text" is XeTeX-only and will crash LuaTeX.
%\font\tenrm="\mainfont:mapping=tex-text" at \mainfontsize
%\font\tenit="\mainfont/I:mapping=tex-text" at \mainfontsize
%\font\tensl="\mainfont:slant=.15;mapping=tex-text" at \mainfontsize
%\font\tenbf="\mainfont/B:mapping=tex-text" at \mainfontsize
%\font\tensc="\mainfont:mapping=tex-text;+smcp" at \mainfontsize
%\font\tentt="\monofont" at \mainfontsize

% Page layout and text-mode extras — not math-related.
%\topskip=\baselineskip \splittopskip=\baselineskip
%\def\makeheadline{...}
%\chardef\ss="DF  \chardef\aa="E5  \chardef\AA="C5  % etc.
% ============== END OF REST OF MATH SETUP ==============
%%%%%%%%%%%%%%%%%%%%%%%%%%%%%%%%%%%%%%%%%%%%%%%%%%%%%%%%%%%%%%%%%%%%%%%%%%%%%%%%
%%%%%%%%%%%%%%%%%%%%%%%%%%%%%%%%%%%%%%%%%%%%%%%%%%%%%%%%%%%%%%%%%%%%%%%%%%%%%%%%
%%%%%%%%%%%%%%%%%%%%%%%%%%%%%%%%%%%%%%%%%%%%%%%%%%%%%%%%%%%%%%%%%%%%%%%%%%%%%%%%

%%%%%%%%%%%%%%%%%%%%%%%%%%%%%%%%%%%%%%%%%%%%%%%%%%%%%%%%%%%%%%%%%%%%%%%%%%%%%%%%
%%%%%%%%%%%%%%%%% comment: Page Direction and Basic Macros %%%%%%%%%%%%%%%%%%%%%
%%%%%%%%%%%%%%%%%%%%%%%%%%%%%%%%%%%%%%%%%%%%%%%%%%%%%%%%%%%%%%%%%%%%%%%%%%%%%%%%

\pagedir TRT \bodydir TRT \pardir TRT \textdir TRT \mathdir TRT
\amiri
%%%%%%%%%%%%%%%%%%%%%%%%%%%%%%%%%%%%%%%%%%%%%%%%%%%%%%%%%%%%%%%%%%%%%%%%%%%%%%%%%%%%%%%%
%%%%%%%%%%%%%%%%%%%%%%%%%%%%%%%%%%%%%%%%%%%%%%%%%%%%%%%%%%%%%%%%%%%%%%%%%%%%%%%%%%%%%%%%
\def\صفحةجديدة{\vfil\eject}
\def\سطرجديد{\hfil\break}
\def\إنجليزي#1{{\textdir TLT #1}}
%%%%%%%%%%%%%%%%%%%%%%%%%%%%%%%%%%%%%%%%%%%%%%%%%%%%%%%%%%%%%%%%%%%%%%%%%%%%%%%%%%%%%%%%
%%%%%%%%%%%%%%%%%%%%%%%%%%%%%%%%%%%%%%%%%%%%%%%%%%%%%%%%%%%%%%%%%%%%%%%%%%%%%%%%%%%%%%%%
\def\مزدوج{«}
\def\مزدوجة{»}

\def\مربع{%
  \ifmmode
    □&
  \else
    $□$%
  \fi
}
%%%%%%%%%%%%%%%%% comment:  تغيير خط الرياضيات %%%%%%%%%%%%%%%%%%%%%%%%%%%%%%%%%%%%%%%%%
%%%%%%%%%%%%%%%%% comment:  تغيير خط الرياضيات %%%%%%%%%%%%%%%%%%%%%%%%%%%%%%%%%%%%%%%%%
%%%%%%%%%%%%%%%%%%%%%%%%%%%%%%%%%%%%%%%%%%%%%%%%%%%%%%%%%%%%%%%%%%%%%%%%%%%%%%%%%%%%%%%%
\let\الصفحةبلا=\nopagenumbers
%%%%%%%%%%%%%%%%%%%%%%%%%%%%%%%%%%%%%%%%%%%%%%%%%%%%%%%%%%%%%%%%%%%%%%%%%%%%%%%%%%%%%%%%
\def\athe#1{\expandafter\directlua{tex.sprint(number_to_eastern_arabic(\the\csname #1\endcsname))}}
\def\rathe#1{\expandafter\directlua{tex.sprint(reversed_number_to_eastern_arabic(\the\csname #1\endcsname))}}
\def\abjadithe#1{\directlua{tex.sprint(number_to_abjadi(\the\csname #1\endcsname))}}
\let\عربية=\athe
\let\عربي=\rathe
\let\أبجدي=\abjadithe
\let\الترقيم=\pageno
%%%%%%%%%%%%%%%%%%%%%%%%%%%%%%%%%%%%%%%%%%%%%%%%%%%%%%%%%%%%%%%%%%%%%%%%%%%%%%%%%%%%%%%%
\let\الرأس=\headline
\let\الذيل=\footline
%%%%%%%%%%%%%%%%%%%%%%%%%%%%%%%%%%%%%%%%%%%%%%%%%%%%%%%%%%%%%%%%%%%%%%%%%%%%%%%%%%%%%%%%

%%%%%%%%%%%%%%%%%%%%%%%%%%%%%%%%%%%%%%%%%%%%%%%%%%%%%%%%%%%%%%%%%%%%%%%%%%%%%%%%
%%%%%%%%%%%%%%%%% comment: Glue, Spacing, and Box Primitives %%%%%%%%%%%%%%%%%%%
%%%%%%%%%%%%%%%%%%%%%%%%%%%%%%%%%%%%%%%%%%%%%%%%%%%%%%%%%%%%%%%%%%%%%%%%%%%%%%%%

\let\حش=\hss
\let\حشو=\hfil
\let\حشوو=\hfill
\let\حشة=\vss
\let\حشوة=\vfil
\let\حشووة=\vfill
%%%%%%%%%%%%%%%%%%%%%%%%%%%%%%%%%%%%%%%%%%%%%%%%%%%%%%%%%%%%%%%%%%%%%%%%%%%%%%%%%%%%%%%%
\def\صمغ#1{%
  \directlua{%
    local input = '\luaescapestring{\unexpanded{#1}}'
    local converted = arabic_to_western_unit(reverse_digit_sequences(input))
    tex.sprint([=[\string\hglue ]=] .. converted .. ' ')
  }%
}

\def\صمغة#1{%
  \directlua{%
    local input = '\luaescapestring{\unexpanded{#1}}'
    local converted = arabic_to_western_unit(reverse_digit_sequences(input))
    tex.sprint([=[\string\vglue ]=] .. converted .. ' ')
  }%
}
\def\قفز#1{%
  \directlua{%
    local input = '\luaescapestring{\unexpanded{#1}}'
    local converted = arabic_to_western_unit(reverse_digit_sequences(input))
    tex.sprint([=[\string\hskip ]=] .. converted .. ' ')
  }%
}

\def\قفزة#1{%
  \directlua{%
    local input = '\luaescapestring{\unexpanded{#1}}'
    local converted = arabic_to_western_unit(reverse_digit_sequences(input))
    tex.sprint([=[\string\vskip ]=] .. converted .. ' ')
  }%
}


\long\def\ارفع#1#2{%
  \directlua{%
    local input1 = '\luaescapestring{\unexpanded{#1}}'
    local input2 = '\luaescapestring{\unexpanded{#2}}'
    local converted = arabic_to_western_unit(reverse_digit_sequences(input1))
    tex.sprint([=[\raise]=] .. converted .. [=[\hbox{]=] .. input2.. [=[}]=])
  }%
}
\long\def\اخفض#1#2{%
  \directlua{%
    local input1 = '\luaescapestring{\unexpanded{#1}}'
    local input2 = '\luaescapestring{\unexpanded{#2}}'
    local converted = arabic_to_western_unit(reverse_digit_sequences(input1))
    tex.sprint([=[\lower]=] .. converted .. [=[\hbox{]=] .. input2.. [=[}]=])       
  }%
}
%%%%%%%%%%%%%%%%%%%%%%%%%%%%%%%%%%%%%%%%%%%%%%%%%%%%%%%%%%%%%%%%%%%%%%%%%%%%%%%%%%%%%%%%
\let\سطريمين=\leftline
\let\سطروسط=\centerline
\let\سطريسار=\rightline
%%%%%%%%%%%%%%%%%%%%%%%%%%%%%%%%%%%%%%%%%%%%%%%%%%%%%%%%%%%%%%%%%%%%%%%%%%%%%%%%%%%%%%%%
\def\ورقةأ٤{\pagewidth=210mm%
\pageheight=297mm%
\hsize=160mm % Width of the text block
\vsize=250mm % Height of the text block
}
%%%%%%%%%%%%%%%%%%%%%%%%%%%%%%%%%%%%%%%%%%%%%%%%%%%%%%%%%%%%%%%%%%%%%%%%%%%%%%%%%%%%%%%%
\let\ملف=\input
%%%%%%%%%%%%%%%%%%%%%%%%%%%%%%%%%%%%%%%%%%%%%%%%%%%%%%%%%%%%%%%%%%%%%%%%%%%%%%%%%%%%%%%%
\let\تباعدالسطور=\baselineskip
%%%%%%%%%%%%%%%%%%%%%%%%%%%%%%%%%%%%%%%%%%%%%%%%%%%%%%%%%%%%%%%%%%%%%%%%%%%%%%%%%%%%%%%%
\let\بلابادئة=\noindent
\let\فقرة=\par
\def\بادئةفقرة{\kern\parindent}
\let\عنصر=\item
\let\عنصرعنصر=\itemitem
\let\رصاصة=\bullet
%%%%%%%%%%%%%%%%% comment:  تغيير خط الرياضيات %%%%%%%%%%%%%%%%%%%%%%%%%%%%%%%%%%%%%%%%%
%%%%%%%%%%%%%%%%% comment:  تغيير خط الرياضيات %%%%%%%%%%%%%%%%%%%%%%%%%%%%%%%%%%%%%%%%%

%%%%%%%%%%%%%%%%%%%%%%%%%%%%%%%%%%%%%%%%%%%%%%%%%%%%%%%%%%%%%%%%%%%%%%%%%%%%%%%%%%%%%%%%
%%%%%%%%%%%%%%%%%%%%% comment: lain LuaTeX Counter Management System %%%%%%%%%%%%%%%%%%%%%%%%%%%%%%%%%%%%%
%%%%%%%%%%%%%%%%%%%%%%%%%%%%%%%%%%%%%%%%%%%%%%%%%%%%%%%%%%%%%%%%%%%%%%%%%%%%%%%%%%%%%%%%
\def\اجعل#1#2{%
  \directlua{%
    local input1 = '\luaescapestring{\unexpanded{#1}}'
    local input2 = '\luaescapestring{\unexpanded{#2}}'
    local converted = arabic_to_western_number(arabic_to_western_unit(reverse_digit_sequences(input2)))
    tex.sprint(input1 .. ' = ' .. converted .. ' ')
  }%
}
%%%%%%%%%%%%%%%%%
\def\newcounter#1{%
	\directlua{%
	 current_counter_name = '\luaescapestring{\unexpanded{#1}}'
	}
  \futurelet\newcounternext\newcountercheck%
}
\def\newcountercheck{%
  \ifx\newcounternext[%
    \expandafter\newcounterwith%
  \else
    \expandafter\newcounterwithout%
  \fi%
}
\def\newcounterwith[#1]{
	\directlua{%
	 	current_counter_parent = '\luaescapestring{\unexpanded{#1}}'
	 	create_counter(current_counter_name, current_counter_parent)
    }
}
\def\newcounterwithout{
	\directlua{%
		create_counter(current_counter_name, "")
    }
}
\let\عدادجديد=\newcounter
%%%%%%%%%%%%%%%%%
\def\resetcounters#1{%
	\directlua{%
		reset_dependent_counters('\luaescapestring{\unexpanded{#1}}')
	}
}
%%%%%%%%%%%%%%%%%
\def\newstatement#1#2#3{%
    \directlua{%
    	local input1 = '\luaescapestring{\unexpanded{#1}}'
    	local input2 = '\luaescapestring{\unexpanded{#2}}'
    	local input3 = '\luaescapestring{\unexpanded{#3}}'
    	create_statement(input1, input2, input3)
    }
}

\def\statement#1{%
	\directlua{%
	current_counter_name = '\luaescapestring{\unexpanded{#1}}'
	increment_statement_counter(current_counter_name)
	}
	\def\statementcontent{\directlua{tex.sprint(print_statement(current_counter_name))}}%
	\bigskip\noindent{\amiribf \hypertarget{\statementcontent}{\statementcontent}}%
}
\let\عبارةجديدة=\newstatement
\let\عبارة=\statement
%%%%%%%%%%%%%%%%%%%%%%%%%%%%%%%%%%%%%%%%%%%%%%%%%%%%%%%%%%%%%%%%%%%%%%%%%%%%%%%%%%%%%%%%
%%%%%%%%%%%%%%%%%%%%% comment: lain LuaTeX Counter Management System %%%%%%%%%%%%%%%%%%%%%%%%%%%%%%%%%%%%%
%%%%%%%%%%%%%%%%%%%%%%%%%%%%%%%%%%%%%%%%%%%%%%%%%%%%%%%%%%%%%%%%%%%%%%%%%%%%%%%%%%%%%%%%




%%%%%%%%%%%%%%%%% comment: التقسيمات %%%%%%%%%%%%%%%%%%%%%%%%%%%%%%%%%%%%%%%%%%%%%%
%%%%%%%%%%%%%%%%% comment: التقسيمات %%%%%%%%%%%%%%%%%%%%%%%%%%%%%%%%%%%%%%%%%%%%%%
% ensure at least #1 vertical space remains, otherwise start new page
\def\ensurevspace#1{%
  \dimen0=\pagegoal \advance\dimen0 by -\pagetotal
  \ifdim\dimen0<#1 \vfill\eject \fi
}


% Counter definitions
\newcount\chaptercounter
\newcount\sectioncounter
\newcount\subsectioncounter
\newcount\bandcounter
\newcount\subsubsectioncounter
% Initialize counters
\chaptercounter=0
\sectioncounter=0
\subsectioncounter=0
\bandcounter=0
\subsubsectioncounter=0



\def\mychapter#1#2{%
  
  }



% Chapter command - resets section and subsection counters
\def\chaptername{\كبيرجدااا الباب}
\def\chapter#1#2{%
\ensurevspace{6\baselineskip}
  \advance\chaptercounter by 1
  \sectioncounter=0
  \subsectioncounter=0
  %\subsubsectioncounter=0
  \resetcounters{باب}
    \vskip 2em plus 1em minus 1em
    %%%%%%%%%%%%%%%%%
    \def\tempa{#1}%
    \def\tempc{و}\def\templ{ش}\def\tempr{ي}%
  \ifx\tempa\tempc
    \centerline{\chaptername\ \rathe{chaptercounter} \hypertarget{\rathe{chaptercounter}}{#2}}%
  \else\ifx\tempa\templ
    \rightline{\chaptername\ \rathe{chaptercounter} \hypertarget{\rathe{chaptercounter}}{#2}}%
  \else\ifx\tempa\tempr
  \leftline{\chaptername\ \rathe{chaptercounter} \hypertarget{\rathe{chaptercounter}}{#2}}%
  \else
    \errmessage{mychapter: unknown alignment '#1', use c, l, or r}%
  \fi\fi\fi
   %%%%%%%%%%%%%%%%%
    \vskip 1.5em plus 0.5em minus 0.5em
  	\noindent
  	
  	  \directlua{
  local input1 =[=[\rathe{chaptercounter}\ \ {\amiribf\quad]=]   .. "\luaescapestring{\string\hyperlink{\rathe{chaptercounter}}{#2}}" .. [=[}\hfill]=]
  local input2 =[=[\rathe{pageno}]=]
  toc_write(input1, input2) 
  }%
}

% Section command - resets subsection counter
\def\sectionname{\كبيرجدا الفصل}
\def\section#1#2{%
\ensurevspace{6\baselineskip}
  \advance\sectioncounter by 1
	\resetcounters{فصل}
  \subsectioncounter=0
  %\subsubsectioncounter=0
    \vskip 1.5em plus 0.5em minus 0.5em
%%%%%%%%%%%%%%%%%
\def\tempa{#1}%
\def\tempc{و}\def\templ{ش}\def\tempr{ي}%
\ifx\tempa\tempc
\centerline{\sectionname\ \rathe{chaptercounter}-\rathe{sectioncounter} \hypertarget{\rathe{chaptercounter}-\rathe{sectioncounter}}{#2}}%
\else\ifx\tempa\templ
\rightline{\sectionname\ \rathe{chaptercounter}-\rathe{sectioncounter} \hypertarget{\rathe{chaptercounter}-\rathe{sectioncounter}}{#2}}%
\else\ifx\tempa\tempr
\leftline{\sectionname\ \rathe{chaptercounter}-\rathe{sectioncounter} \hypertarget{\rathe{chaptercounter}-\rathe{sectioncounter}}{#2}}%
\else
\errmessage{mychapter: unknown alignment '#1', use c, l, or r}%
\fi\fi\fi
%%%%%%%%%%%%%%%%%
    \vskip 1em plus 0.3em minus 0.3em
  	\noindent
  	
  	  \directlua{
  local input1 =[=[\quad\quad\rathe{chaptercounter}-\rathe{sectioncounter}\quad]=]   .. "\luaescapestring{\string\hyperlink{\rathe{chaptercounter}-\rathe{sectioncounter}}{#2}}" .. "\string\\dotfill"
  local input2 =[=[\rathe{pageno}]=]
  toc_write(input1, input2) 
  }%
}

% Subsection command
\def\subsectionname{\كبير المبحث}
\def\subsection#1#2{%
  \advance\subsectioncounter by 1
  \resetcounters{مبحث}
  %\subsubsectioncounter=0
    \vskip 1em plus 0.3em minus 0.3em
    %%%%%%%%%%%%%%%%%
    %%%%%%%%%%%%%%%%%
\def\tempa{#1}%
\def\tempc{و}\def\templ{ش}\def\tempr{ي}%
\ifx\tempa\tempc
\centerline{\subsectionname \rathe{chaptercounter}-\rathe{sectioncounter}-\rathe{subsectioncounter} \hypertarget{\rathe{chaptercounter}-\rathe{sectioncounter}-\rathe{subsectioncounter}}{#2}}%
\else\ifx\tempa\templ
\rightline{\subsectionname \rathe{chaptercounter}-\rathe{sectioncounter}-\rathe{subsectioncounter} \hypertarget{\rathe{chaptercounter}-\rathe{sectioncounter}-\rathe{subsectioncounter}}{#2}}%
\else\ifx\tempa\tempr
\leftline{\subsectionname \rathe{chaptercounter}-\rathe{sectioncounter}-\rathe{subsectioncounter} \hypertarget{\rathe{chaptercounter}-\rathe{sectioncounter}-\rathe{subsectioncounter}}{#2}}%
\else
\errmessage{mychapter: unknown alignment '#1', use c, l, or r}%
\fi\fi\fi
%%%%%%%%%%%%%%%%%
  	\vskip 0.5em plus 0.2em minus 0.2em
  	\noindent
  	
  	  	  \directlua{
  local input1 =[=[\ \quad\quad\quad\quad \rathe{chaptercounter}-\rathe{sectioncounter}-\rathe{subsectioncounter}\quad]=]   .. "\luaescapestring{\string\hyperlink{\rathe{chaptercounter}-\rathe{sectioncounter}-\rathe{subsectioncounter}}{#2}}" .. "\string\\dotfill"
  local input2 =[=[\rathe{pageno}]=]
  toc_write(input1, input2) 
  }%
}

% Subsection command
%\def\subsectionname{المبحث}
%\def\subsection{%
  %\advance\subsubsectioncounter by 1
    %\par  % ensure we are in vertical mode
  %\vskip 1.5\baselineskip
  %\centerline{* * *}%
  %\vskip 1.5\baselineskip
 %\par{\amiribf\subsectionname\ \hypertarget{\abjadithe{chaptercounter}-\abjadithe{subsubsectioncounter}}{\rathe{subsubsectioncounter}}.}
%}

\newcount\savedchaptercounter
\savedchaptercounter=\chaptercounter
\newcount\savedsectioncounter
\savedsectioncounter=\sectioncounter
\def\band#1{%
  \advance\bandcounter by 1
  \ifnum\bandcounter>1
    % Not the first subsection: check if chapter/section unchanged
    \ifnum\savedchaptercounter=\chaptercounter
      \ifnum\savedsectioncounter=\sectioncounter
        % Both counters unchanged → print separator
%        \par
%        \vskip 0.3\baselineskip
%        \centerline{* * *}%
%        \vskip 0.3\baselineskip
\bigskip
      \else
        % Section changed: update saved counters, no separator
        \savedsectioncounter=\sectioncounter
        \savedchaptercounter=\chaptercounter
      \fi
    \else
      % Chapter changed: update saved counters, no separator
      \savedsectioncounter=\sectioncounter
      \savedchaptercounter=\chaptercounter
    \fi
  \else
    % First subsection: update saved counters to current values
    \savedsectioncounter=\sectioncounter
    \savedchaptercounter=\chaptercounter
  \fi
  \par
  {\amiribf#1 \hypertarget{#1 \rathe{bandcounter}}{\rathe{bandcounter}}. }%
}
%%%%%%%%%%%%%%%%%
\let\باب=\chapter
\let\عدادالباب=\chaptercounter
\let\فصل=\section
\let\عدادالفصل=\sectioncounter
\let\مبحث=\subsection
\let\عدادالمبحث=\subsectioncounter
\def\بند{\band{البند}}
\let\عدادالبند=\bandcounter
%%%%%%%%%%%%%%%%% comment: التقسيمات %%%%%%%%%%%%%%%%%%%%%%%%%%%%%%%%%%%%%%%%%%%%%%
%%%%%%%%%%%%%%%%% comment: التقسيمات %%%%%%%%%%%%%%%%%%%%%%%%%%%%%%%%%%%%%%%%%%%%%%
\newcount\figurecounter
\figurecounter=0
\def\figure{%
  \global\advance\figurecounter by 1%
  \def\currentcontent{الشكل \rathe{figurecounter}}%
  \hypertarget{\currentcontent}{\currentcontent}%
}
\let\شكل=\figure


\newcount\propositioncounter
\propositioncounter=0
\def\proposition#1#2{%
  \global\advance\propositioncounter by 1%
  \hypertarget{#1 \rathe{propositioncounter} #2}{\bigskip\noindent{\amiribi #1 \rathe{propositioncounter} (#2):}}%
}
\let\قضية=\proposition

\newcount\definitioncounter
\definitioncounter=0
\def\definition#1{%
  \global\advance\definitioncounter by 1%
  \hypertarget{تعريف \rathe{definitioncounter} #1}{\bigskip\noindent{\amiribi تعريف \rathe{definitioncounter} #1:}}%
}
\let\تعريف=\definition


\newcount\examplecounter
\examplecounter=0
\def\example{%
  \global\advance\examplecounter by 1%
  \hypertarget{مثال \rathe{examplecounter}}{\bigskip\noindent{\amiribi مثال \rathe{examplecounter}: }}%
}
\let\مثال=\example


\newcount\axiomcounter
\axiomcounter=0
\def\axiom{%
  \global\advance\axiomcounter by 1%
  \hypertarget{مسلَّمة \rathe{axiomcounter}}{\bigskip\noindent{\amiribi مسلَّمة \rathe{axiomcounter}: }}%
}
\let\مسلمة=\example
%%%%%%%%%%%%%%%%%%%%%%%%%%%%%%%%%%%%%%%%%%%%%%%%%%%%%%%%%%%%%%%%%%%%%%%%%%%%%%%%
%%%%%%%%%%%%%%%%% comment: Footnotes %%%%%%%%%%%%%%%%%%%%%%%%%%%%%%%%%%%%%%%%%%%
%%%%%%%%%%%%%%%%%%%%%%%%%%%%%%%%%%%%%%%%%%%%%%%%%%%%%%%%%%%%%%%%%%%%%%%%%%%%%%%%

%%%%%%%%%%%%%%%%% comment:  الحاشية %%%%%%%%%%%%%%%%%%%%%%%%%%%%%%%%%%%%%%%%%%%%%%%
%%%%%%%%%%%%%%%%% comment:  الحاشية %%%%%%%%%%%%%%%%%%%%%%%%%%%%%%%%%%%%%%%%%%%%%%%
%%%%%%%%%%%%%%%%%
\newcount\fnabscount
\fnabscount=0
\def\autofootnote#1{%
    \global\advance\fnabscount by 1%
    \edef\thisabs{\the\fnabscount}%
    \latelua{record_fn_page(tex.count[0])}%
    \edef\currentfootnum{\directlua{tex.sprint(reversed_number_to_eastern_arabic(get_footnote_num(\thisabs)))}}%
    \footnote{\hyperlink{\thisabs}{${}^{\currentfootnum}$}}{\hypertarget{\thisabs}{#1}}%
}
\let\حاشية=\autofootnote
%%%%%%%%%%%%%%%%%
\let\الحاشية=\footnote
%%%%%comment: مثال على الاستخدام \الحاشية{${}^١$}{نص الحاشية}
\let\حاشية=\autofootnote
%%%%%comment: مثال على الاستخدام \حاشية{نص الحاشية}
\let\عدادالحاشية=\footnotecounter
%%%%%%%%%%%%%%%%% comment:  الحاشية %%%%%%%%%%%%%%%%%%%%%%%%%%%%%%%%%%%%%%%%%%%%%%%
%%%%%%%%%%%%%%%%% comment:  الحاشية %%%%%%%%%%%%%%%%%%%%%%%%%%%%%%%%%%%%%%%%%%%%%%%

\def\مجوف#1{
\pdfextension literal{1 Tr 0.5 w}%
#1
\pdfextension literal{0 Tr}%
}

\def\مسمك#1{%
  \pdfextension literal{2 Tr 0.3 w}%
  #1%
  \pdfextension literal{0 Tr}%
}





\newbox\italboxA
\newbox\italboxB

% #1 = slant value (e.g. 0.22 forward, -0.22 backslant)
% #2 = content
\def\fauxitalic#1#2{%
  \leavevmode
  %
  % Measure the original content
  \setbox\italboxA\hbox{{#2}}%
  %
  % Apply PDF shear in a zero-width box
  \setbox\italboxB\hbox{%
    \pdfextension literal{q 1 0 #1 1 0 0 cm}%
    \rlap{\copy\italboxA}%
    \pdfextension literal{Q}%
  }%
  %
  % Preserve original height/depth (slant doesn't change them)
  \ht\italboxB\ht\italboxA
  \dp\italboxB\dp\italboxA
  %
  % Overhang = |slant| × height
  \ifdim #1pt<0pt
    % Backslant: top bleeds LEFT, so kern right before placing
    \hbox to\dimexpr\wd\italboxA + -#1\ht\italboxA\relax{%
      \kern\dimexpr-#1\ht\italboxA\relax   % push start point right
      \box\italboxB
      \hss
    }%
  \else
    % Forward slant: top bleeds RIGHT, extra space on right
    \hbox to\dimexpr\wd\italboxA + #1\ht\italboxA\relax{%
      \box\italboxB
      \hss
    }%
  \fi
}

\def\مميل#1{\fauxitalic{-0.22}{$#1$}}

%%%%%%%%%%%%%%%%%%%%%%%%%%%%%%%%%%%%%%%%%%%%%%%%%%%%%%%%%%%%%%%%%%%%%%%%%%%%%%%%
%%%%%%%%%%%%%%%%% comment: Strut and Padding for Boxes %%%%%%%%%%%%%%%%%%%%%%%%%
%%%%%%%%%%%%%%%%%%%%%%%%%%%%%%%%%%%%%%%%%%%%%%%%%%%%%%%%%%%%%%%%%%%%%%%%%%%%%%%%

%%%%%%%%%%%%%%%%% comment:  دعامة %%%%%%%%%%%%%%%%%%%%%%%%%%%%%%%%%%%%%%%%%%%%%%
%%%%%%%%%%%%%%%%% comment:  دعامة %%%%%%%%%%%%%%%%%%%%%%%%%%%%%%%%%%%%%%%%%%%%%%
\newbox\mystrutbox
\def\computestrut{%
  \setbox\mystrutbox=\hbox{%
    \vrule height 0.7\baselineskip depth 0.3\baselineskip width 0pt%
  }%
}
\def\mystrut{\copy\mystrutbox}
\let\دعامة=\mystrut

\newbox\padbox
\newif\ifpadding
\long\def\padcell#1{\setbox\padbox=\hbox{#1}%
  \paddingtrue
  \ifdim\ht\padbox=0pt \ifdim\dp\padbox=0pt
    \paddingfalse
  \fi\fi
  \ifpadding
    \ht\padbox=\dimexpr\ht\padbox+4pt\relax
    \dp\padbox=\dimexpr\dp\padbox+4pt\relax
  \fi
  \box\padbox}
  
  \let\خلية=\padcell
%%%%%%%%%%%%%%%%% comment:  دعامة %%%%%%%%%%%%%%%%%%%%%%%%%%%%%%%%%%%%%%%%%%%%%%
%%%%%%%%%%%%%%%%% comment:  دعامة %%%%%%%%%%%%%%%%%%%%%%%%%%%%%%%%%%%%%%%%%%%%%%

%%%%%%%%%%%%%%%%%%%%%%%%%%%%%%%%%%%%%%%%%%%%%%%%%%%%%%%%%%%%%%%%%%%%%%%%%%%%%%%%
%%%%%%%%%%%%%%%%% comment: Math Symbols and Operators %%%%%%%%%%%%%%%%%%%%%%%%%%
%%%%%%%%%%%%%%%%%%%%%%%%%%%%%%%%%%%%%%%%%%%%%%%%%%%%%%%%%%%%%%%%%%%%%%%%%%%%%%%%

%%%%%%%%%%%%%%%%% comment:  رياضيات %%%%%%%%%%%%%%%%%%%%%%%%%%%%%%%%%%%%%%%%%%%%%%
%%%%%%%%%%%%%%%%% comment:  رياضيات %%%%%%%%%%%%%%%%%%%%%%%%%%%%%%%%%%%%%%%%%%%%%%
% Safe dot macros (no internal registers)
\def\ldots{%
  \mathinner{%
    \hbox{$.$}\hbox{$.$}\hbox{$.$}%
  }%
} % horizontal low dots (baseline dots)

\def\cdots{%
  \mathinner{%
    \vcenter{\hbox{\hbox{$.$}\hbox{$.$}\hbox{$.$}}}%
  }%
} % centered horizontal dots

\def\vdots{%
  \mathinner{%
    \vcenter{\baselineskip4pt \lineskiplimit0pt
      \kern6pt \hbox{$.$}\hbox{$.$}\hbox{$.$}}%
  }%
} % vertical dots

% ddots: diagonal dots, scale a bit depending on math style
\def\ddots{%
  \mathinner{%
    \mathchoice
      {\mkern1mu\raise7pt\hbox{$.$}\mkern2mu\raise2pt\hbox{$.$}\mkern2mu\lower3pt\hbox{$.$}}% \displaystyle
      {\mkern1mu\raise5pt\hbox{$.$}\mkern2mu\raise1.5pt\hbox{$.$}\mkern2mu\lower2.5pt\hbox{$.$}}% \textstyle
      {\mkern1mu\raise4pt\hbox{$.$}\mkern2mu\raise1pt\hbox{$.$}\mkern2mu\lower1.8pt\hbox{$.$}}% \scriptstyle
      {\mkern1mu\raise3pt\hbox{$.$}\mkern2mu\raise0.5pt\hbox{$.$}\mkern2mu\lower1.2pt\hbox{$.$}}% \scriptscriptstyle
  }%
}

\def\lddots{%
  \mathinner{%
    \mathchoice
      {\mkern1mu\lower3pt\hbox{$.$}\mkern2mu\raise2pt\hbox{$.$}\mkern2mu\raise7pt\hbox{$.$}}% \displaystyle (reverse of \حذفةقطرية)
      {\mkern1mu\lower2.5pt\hbox{$.$}\mkern2mu\raise1.5pt\hbox{$.$}\mkern2mu\raise5pt\hbox{$.$}}% \textstyle
      {\mkern1mu\lower1.8pt\hbox{$.$}\mkern2mu\raise1pt\hbox{$.$}\mkern2mu\raise4pt\hbox{$.$}}% \scriptstyle
      {\mkern1mu\lower1.2pt\hbox{$.$}\mkern2mu\raise0.5pt\hbox{$.$}\mkern2mu\raise3pt\hbox{$.$}}% \scriptscriptstyle
  }%
}

\def\frac#1#2{%
	\mathinner{
		\mathchoice
		{\vcenter{\hbox{$\displaystyle #1$}}\kern-.2ex\over\vcenter{\hbox{$\displaystyle #2$}}}%
		{\vcenter{\hbox{$\textstyle #1$}}\kern-.2ex\over\vcenter{\hbox{$\textstyle #2$}}}%
		{\vcenter{\hbox{$\scriptstyle #1$}}\kern-.15ex\over\vcenter{\hbox{$\scriptstyle #2$}}}%
		{\vcenter{\hbox{$\scriptscriptstyle #1$}}\kern-.1ex\over\vcenter{\hbox{$\scriptscriptstyle #2$}}}%
	}
}

\def\symmtric#1{%
	\mathinner{
		\mathchoice
		{\vcenter{\hbox{$\displaystyle #1$}}}%
		{\vcenter{\hbox{$\textstyle #1$}}}%
		{\vcenter{\hbox{$\scriptstyle #1$}}}%
		{\vcenter{\hbox{$\scriptscriptstyle #1$}}}%
	}
}
%%%%%%%%%%%%%%%%%%%%%%%%%%%%%%%%%%%%%%%%%%%%%%%%%%%%%%%%%%%%%%%%%%%%%%%%%%%%%%%%%%%%%%%%
\def\arabicsqrt#1#2{%
  \mathinner{\reflect{$\textdir TLT#1{\sqrt{\reflect{$\textdir TRT#1 #2$}}}$}}%
}

\def\asqrt#1{%
  \mathchoice
   {\arabicsqrt{\displaystyle}{#1}}%
    {\arabicsqrt{\textstyle}{#1}}%
    {\arabicsqrt{\scriptstyle}{#1}}%
    {\arabicsqrt{\scriptscriptstyle}{#1}}%
}
\def\arabicroot#1#2#3#4{%
  \mathinner{\reflect{$\mathdir TLT \textdir TLT{#1\root{\reflect{$\mathdir TRT \textdir TRT #2 {#3}$}}\of{\reflect{$\mathdir TRT \textdir TRT #1 {#4}$}}}$}}
}

\def\aroot#1#2{%
		\mathchoice
		{\arabicroot{\displaystyle}{\scriptstyle}{#1}{#2}}%
		{\arabicroot{\textstyle}{\scriptstyle}{#1}{#2}}%
		{\arabicroot{\scriptstyle}{\scriptscriptstyle}{#1}{#2}}%
		{\arabicroot{\scriptscriptstyle}{\scriptscriptstyle}{#1}{#2}}%
}


\def\fullaroot{\futurelet\fullarootnext\fullarootcheck}
\def\fullarootcheck{%
  \ifx\fullarootnext[%           % is the next token an open bracket?
    \let\fullarootnext\fullarootwith%   yes → switch to the width-aware handler
  \else
    \let\fullarootnext\fullarootwithout% no  → switch to the plain handler
  \fi
  \fullarootnext}                % call whichever handler was selected
\def\fullarootwith[#1]#2{\aroot{#1}{#2}}
\def\fullarootwithout#1{\asqrt{#1}}
\let\جذر=\fullaroot

%%%%%%%%%%%%%%%%%%%%%%%%%%%%%%%%%%%%%%%%%%%%%%%%%%%%%%%%%%%%%%%%%%%%%%%%%%%%%%%%%%%%%%%%


\def\doublesignshelper#1#2#3{%
  \mathbin{%
    \vcenter{%
      \offinterlineskip
      \halign{\hfil$#1##$\hfil\cr
        #2\cr
        #3\cr
      }%
    }%
  }%
}

\def\doublesigns#1#2{%
  \mathchoice
    {\doublesignshelper{\displaystyle}{#1}{#2}}%
    {\doublesignshelper{\textstyle}{#1}{#2}}%
    {\doublesignshelper{\scriptstyle}{#1}{#2}}%
    {\doublesignshelper{\scriptscriptstyle}{#1}{#2}}%
}


\def\hidewidth{\hskip -1000pt plus 1fill}
\let\بلاارتفاع=\smash
\def\clap#1{\hbox to 0pt{\hss#1\hss}}
\def\overlap#1#2{%
\def\tempa{#1}%
\def\tempc{و}\def\templ{ش}\def\tempr{ي}%
\ifx\tempa\tempc
\clap{#2}%
\else\ifx\tempa\templ
\llap{#2}%
\else\ifx\tempa\tempr
\rlap{#2}%
\else
\errmessage{overlap: unknown alignment '#1', use c, l, or r}%
\fi\fi\fi
}
\let\تراكب=\overlap
\newbox\mystrutbox
\def\computestrut{%
  \setbox\mystrutbox=\hbox{%
    \vrule height 0.7\baselineskip depth 0.3\baselineskip width 0pt%
  }%
}
\def\mystrut{\copy\mystrutbox}


\long\def\equation#1#2{%
\mathinner{%
\mathchoice
	{\directlua{%
    		tex.sprint(process_equations ('\luaescapestring{\unexpanded{\displaystyle}}', '\luaescapestring{\unexpanded{#1}}', '\luaescapestring{\unexpanded{#2}}'))
  }}%
	{\directlua{%
    		tex.sprint(process_equations ('\luaescapestring{\unexpanded{\textstyle}}', '\luaescapestring{\unexpanded{#1}}', '\luaescapestring{\unexpanded{#2}}'))
  }}%
	{\directlua{%
    		tex.sprint(process_equations ('\luaescapestring{\unexpanded{\scriptstyle}}', '\luaescapestring{\unexpanded{#1}}', '\luaescapestring{\unexpanded{#2}}'))
  }}%
	{\directlua{%
    		tex.sprint(process_equations ('\luaescapestring{\unexpanded{\scriptscriptstyle}}', '\luaescapestring{\unexpanded{#1}}', '\luaescapestring{\unexpanded{#2}}'))
  }}%
}%
}
\long\def\numberedequation#1#2{%
\mathinner{%
\mathchoice
	{\directlua{%
    		tex.sprint(process_numbered_equations ('\luaescapestring{\unexpanded{\displaystyle}}', '\luaescapestring{\unexpanded{#1}}', '\luaescapestring{\unexpanded{#2}}'))
  }}%
	{\directlua{%
    		tex.sprint(process_numbered_equations ('\luaescapestring{\unexpanded{\textstyle}}', '\luaescapestring{\unexpanded{#1}}', '\luaescapestring{\unexpanded{#2}}'))
  }}%
	{\directlua{%
    		tex.sprint(process_numbered_equations ('\luaescapestring{\unexpanded{\scriptstyle}}', '\luaescapestring{\unexpanded{#1}}', '\luaescapestring{\unexpanded{#2}}'))
  }}%
	{\directlua{%
    		tex.sprint(process_numbered_equations ('\luaescapestring{\unexpanded{\scriptscriptstyle}}', '\luaescapestring{\unexpanded{#1}}', '\luaescapestring{\unexpanded{#2}}'))
  }}%
}%
}

\long\def\widetable#1#2{%
	\directlua{%
    		tex.sprint(process_wide_table ('\luaescapestring{\unexpanded{#1}}', '\luaescapestring{\unexpanded{#2}}'))
  }%
}%

\long\def\table#1#2{%
	\directlua{%
    		tex.sprint(process_table ('\luaescapestring{\unexpanded{#1}}', '\luaescapestring{\unexpanded{#2}}'))
  }%
}%

\def\تقريب{%
	\mathrel{
		\mathchoice
		{\reflect{$\displaystyle\approx$}}%
		{\reflect{$\textstyle\approx$}}%
		{\reflect{$\scriptstyle\approx$}}%
		{\reflect{$\scriptscriptstyle\approx$}}%
	}
}
\def\فرق{%
	\mathbin{
		\mathchoice
		{\reflect{$\displaystyle \setminus$}}%
		{\reflect{$\textstyle \setminus$}}%
		{\reflect{$\scriptstyle \setminus$}}%
		{\reflect{$\scriptscriptstyle \setminus$}}%
	}
}


\def\واصلة#1{%
  \ifmmode
	\overline{#1}
  \else
    $\overline{\hbox{#1}}$%
  \fi
}
\def\واصل#1{%
  \ifmmode
	\underline{#1}
  \else
    $\underline{\hbox{#1}}$%
  \fi
}


\def\coloneqq{\mathrel{%
  \mathopen{\mathchoice
    {\raise.20ex\hbox{$\displaystyle:$}}%  display
    {\raise.21ex\hbox{$\textstyle:$}}%  text
    {\raise.10ex\hbox{$\scriptstyle:$}}%  script
    {\raise.05ex\hbox{$\scriptscriptstyle:$}}%  scriptscript
  }%
  \mkern-1.2mu=}}
  
\def\afactorial#1{%
	\mathinner{%
		\mathchoice
		{\vrule\vbox{\kern2pt\hbox{\kern2pt$\displaystyle#1$\kern2pt}\kern2pt\hrule}}%
		{\vrule\vbox{\kern2pt\hbox{\kern2pt$\textstyle#1$\kern2pt}\kern2pt\hrule}}%
		{\vrule\vbox{\kern2pt\hbox{\kern2pt$\scriptstyle#1$\kern2pt}\kern2pt\hrule}}%
		{\vrule\vbox{\kern2pt\hbox{\kern2pt$\scriptscriptstyle#1$\kern2pt}\kern2pt\hrule}}%
	}%
}

\def\closedintegral#1{%
  \def\closedintegralarg{#1}%        % save in a macro, not in input stream!
  \futurelet\closedintegralnext\closedintegralcheck%
}
\def\closedintegralcheck{%
  \ifx\closedintegralnext[%
    \expandafter\closedintegralwith%
  \else
    \expandafter\closedintegralwithout%
  \fi%
}
\def\closedintegralwith[#1]{\oint_{\closedintegralarg}^{#1}}
\def\closedintegralwithout{\oint_{\closedintegralarg}}


\def\معقوفةفوق#1{%
    \mathop{%
        \mathchoice%
        {\reflect{\hbox{\mathdir TLT$\displaystyle{\overbrace{\reflect{\mathdir TRT$\displaystyle#1$}}}$}}}%
        {\reflect{\hbox{\mathdir TLT$\textstyle{\overbrace{\reflect{\mathdir TRT$\textstyle#1$}}}$}}}%
        {\reflect{\hbox{\mathdir TLT$\scriptstyle{\overbrace{\reflect{\mathdir TRT$\scriptstyle#1$}}}$}}}%
        {\reflect{\hbox{\mathdir TLT$\scriptscriptstyle{\overbrace{\reflect{\mathdir TRT$\scriptscriptstyle#1$}}}$}}}%
    }
}

\def\معقوفةتحت#1{%
    \mathop{%
        \mathchoice%
        {\reflect{\hbox{\mathdir TLT$\displaystyle{\underbrace{\reflect{\mathdir TRT$\displaystyle#1$}}}$}}}%
        {\reflect{\hbox{\mathdir TLT$\textstyle{\underbrace{\reflect{\mathdir TRT$\textstyle#1$}}}$}}}%
        {\reflect{\hbox{\mathdir TLT$\scriptstyle{\underbrace{\reflect{\mathdir TRT$\scriptstyle#1$}}}$}}}%
        {\reflect{\hbox{\mathdir TLT$\scriptscriptstyle{\underbrace{\reflect{\mathdir TRT$\scriptscriptstyle#1$}}}$}}}%
    }
}
%%%%%%%%%%%%%%%%%
\let\كنص=\textstyle
\let\كأساس=\displaystyle
\let\كأس=\scriptstyle
\let\كأسأس=\scriptscriptstyle
\let\يمين=\left
\let\يسار=\right
\let\يمينها=\leqno
\let\يسارها=\eqno
\let\تباعدالأعمدة=\tabskip
\let\عد=\cr
\let\احذف=\omit
\let\امدد=\span
\let\حشوخط=\hrulefill
\let\حشونقاط=\dotfill
\def\حاجب{\omit\hrule}
\def\حاجبة{\omit\vrule}
\let\خط=\hrule
\let\خطة=\vrule
\let\قاعدة=\hrule
\let\عمود=\vrule
\let\بلاتنسيق=\noalign
\def\حجاب{\noalign{\hrule}}
\let\شرطة=\prime
\let\حذفةسطرية=\ldots
\let\حذفةأفقية=\cdots
\let\حذفةرأسية=\vdots
\let\حذفةقطرية=\ddots
\let\حذفةقطريةمعكوسة=\lddots
\let\كسر=\frac
\let\مصفوف=\equation
\let\مصفوفة=\matrix
\let\خطوات=\numberedequation
\def\موجبةسالبة{±}
\def\سالبةموجبة{∓}
\def\onbaseline#1{%
  \setbox0=\hbox{#1}%
  \raise-\dp0\box0%
}
\Umathchardef\arabicplus="2"0"2A3C        % ⨼
\Umathchardef\arabicstar="2"0"066D        % ٭
\let\زائد=\arabicplus
\let\ضرب=\arabicstar
\let\في=\cdot
\def\نقطة#1{\mathaccent"02D9 #1}
\let\شرطةكسر=\backslash
\let\إشارةمزدوجة=\doublesigns
\let\أعمدة=\widetable
\let\جدول=\table
\let\يطابق=\equiv
\let\يكافئ=\iff
\def\أصغرأويساوي{≥}
\def\أكبرأويساوي{≤}
%% Left arrows (مذكر) — Unicode code points
\def\سهيم{←}	% ←  U+2190
\def\إلى{←}	% ←  U+2190
\def\سهم{⟵}	% ⟵  U+27F5
\def\رميح{⇐}	% ⇐  U+21D0
\def\رمح{⟸}	% ⟸  U+27F8  
%% Right arrows (مؤنث) — Unicode code points
\def\سهيمة{→}% →  U+2192
\def\سهمة{⟶}% ⟶  U+27F6
\def\رميحة{⇒}% ⇒  U+21D2
\def\رمحة{⟹}% ⟹  U+27F9
\let\خيال=\hphantom
\let\خيالة=\vphantom
\let\خيالات=\phantom
\let\يتخذ=\coloneqq
\def\الخالية{∅}
\def\بماأن{∵}
\def\فإن{∴}
\def\جزئية{⊃}
\def\غيرجزئية{⊅}
\def\جزئيةفعلية{⊋}
\def\جزئيةأوتساوي{⊇}
\def\شاملة{⊂}
\def\غيرشاملة{⊄}
\def\شاملةفعلية{⊊}
\def\شاملةأوتساوي{⊆}
\def\ينتمي{∋}
\def\لاينتمي{∌}
\def\يحتوي{∈}
\def\لايحتوي{∉}
\def\لايساوي{≠}
\def\اتحاد{∪}
\def\تقاطع{∩}
\def\اتحادات{⋃}
\def\تقاطعات{⋂}
\def\مالانهاية{∞}
\let\قفزةكبيرة=\bigskip
\let\قفزةمتوسطة=\medskip
\let\قفزةصغيرة=\smallskip
\let\تراتيب=\afactorial
\let\متماثل=\symmtric
\let\دائرة=\circ
\let\تباعدعلاقة=\mathrel
\def\نجمة{⋆}
\def\يشابه{∽}
\def\يشبه{∼}
\def\و{∧}
\def\أو{∨}
\def\ليس{⌐}
\def\لايقتضي{⇍}
\def\لايشترط{⇏}
\let\قاسم=\mid
%%%%%%%%%%%%%%%%% comment:  رياضيات %%%%%%%%%%%%%%%%%%%%%%%%%%%%%%%%%%%%%%%%%%%%%%
%%%%%%%%%%%%%%%%% comment:  رياضيات %%%%%%%%%%%%%%%%%%%%%%%%%%%%%%%%%%%%%%%%%%%%%%

%%%%%%%%%%%%%%%%%%%%%%%%%%%%%%%%%%%%%%%%%%%%%%%%%%%%%%%%%%%%%%%%%%%%%%%%%%%%%%%%
%%%%%%%%%%%%%%%%% comment: Image Handling (Scaling, Rotation, Inclusion) %%%%%%%
%%%%%%%%%%%%%%%%%%%%%%%%%%%%%%%%%%%%%%%%%%%%%%%%%%%%%%%%%%%%%%%%%%%%%%%%%%%%%%%%

%%%%%%%%%%%%%%%%%%%%%%%%%%%%%%%%%%%%%%%%%%%%%%%%%%%%%%%%%%%%%%%%%%%%%%%%
%%%%%%%%%%%%%%%%% comment:  صور %%%%%%%%%%%%%%%%%%%%%%%%%%%%%%%%%%%%%%%%%%%%%%
%%%%%%%%%%%%%%%%%%%%%%%%%%%%%%%%%%%%%%%%%%%%%%%%%%%%%%%%%%%%%%%%%%%%%%%%
% Define scratch registers (equivalent to LaTeX's internal ones)
\newbox\scaleboxA  % Equivalent to box register 0
\newbox\scaleboxB  % Equivalent to box register 2

% Define a scaling box macro with horizontal and vertical scaling factors
% Parameters:
%   #1 = horizontal scale factor (e.g., 0.7 for 70% width)
%   #2 = vertical scale factor (e.g., 1.0 for 100% height)
%   #3 = content to be scaled
\long\def\scalebox#1#2#3{%
  % Ensure we're in horizontal mode
  \leavevmode
  %
  % Store the scale factors in macros
  \def\hscale{#1}%
  \def\vscale{#2}%
  %
  % Put the content in a box so we can measure it
  \setbox\scaleboxA\hbox{{#3}}%
  %
  % Create the scaled box using PDF transformation
  \setbox\scaleboxB\hbox{%
    % Begin PDF transformation group
    % The 'q' saves the graphics state, 'cm' applies transformation matrix
    % Transformation matrix: [sx 0 0 sy 0 0] where sx = horizontal scale, sy = vertical scale
    \pdfextension literal {q \hscale\space 0 0 \vscale\space 0 0 cm}%
    % Overlay the original content (zero-width, positioned relative to current point)
    \rlap{\copy\scaleboxA}%
    % End PDF transformation group (restores graphics state)
    \pdfextension literal {Q}%
  }%
  %
  % Adjust vertical dimensions (height and depth) based on vertical scaling
  % Check if vertical scale factor is negative (flipping vertically)
  \ifdim #2pt<0pt
    % Negative vertical scaling: swap height and depth, make them positive
    % Height of new box = -(vertical scale factor × depth of original box)
    % Depth of new box = -(vertical scale factor × height of original box)
    \ht\scaleboxB-\dimexpr#2\dp\scaleboxA\relax
    \dp\scaleboxB-\dimexpr#2\ht\scaleboxA\relax
  \else
    % Positive vertical scaling: scale height and depth normally
    % Height of new box = vertical scale factor × height of original box
    % Depth of new box = vertical scale factor × depth of original box
    \ht\scaleboxB\dimexpr#2\ht\scaleboxA\relax
    \dp\scaleboxB\dimexpr#2\dp\scaleboxA\relax
  \fi
  %
  % Adjust horizontal dimensions based on horizontal scaling
  % Check if horizontal scale factor is negative (mirroring horizontally)
  \ifdim #1pt<0pt
    % Negative horizontal scaling: create left-aligned box with mirrored content
    % Total width = absolute value of (horizontal scale factor × width of original box)
    \hbox to-\dimexpr#1\wd\scaleboxA\relax{%
      % Shift content to account for mirroring (aligns the right edge)
      \kern-\dimexpr#1\wd\scaleboxA\relax
      % Place the transformed box
      \box\scaleboxB
      % Fill any remaining space with hss (horizontal stretch/shrink)
      \hss
    }%
  \else
    % Positive horizontal scaling: create right-aligned box
    % Total width = horizontal scale factor × width of original box
    \hbox to\dimexpr#1\wd\scaleboxA\relax{%
      % Place the transformed box
      \box\scaleboxB
      % Add kern to fill space (since box might be smaller than target width)
      \kern\dimexpr#1\wd\scaleboxA\relax
      % Fill any remaining space with hss (horizontal stretch/shrink)
      \hss
    }%
  \fi
}
\def\reflect#1{\scalebox{-1}{1}{#1}}

\def\تحجيم#1#2#3{%
  \directlua{%
    tex.sprint(arabic_scale_box('\luaescapestring{\unexpanded{#1}}', '\luaescapestring{\unexpanded{#2}}', '\luaescapestring{\unexpanded{#3}}'))
  }%
}
%%%%%%%%%%%%%%%%%%%%%%%%%%%%%%%%%%%%%%%%%%%%%%%%%%%%%%%%%%%%%%%%%%%%%%%%
%%%%%%%%%%%%%%%%%%%%%%%%%%%%%%%%%%%%%%%%%%%%%%%%%%%%%%%%%%%%%%%%%%%%%%%%
%%%%%%%%%%%%%%%%%%%%%%%%%%%%%%%%%%%%%%%%%%%%%%%%%%%%%%%%%%%%%%%%%%%%%%%%
% ============================================================
% rotatebox + scalebox for LuaTeX plain TeX
% ============================================================
\catcode`@=11

% ---- kernel helpers ----------------------------------------
\def\zap@space#1 #2{#1\ifx#2\@empty\else\expandafter\zap@space\fi#2}
\def\@car#1#2\@nil{#1}
\def\@cdr#1#2\@nil{#2}
\def\@empty{}
\long\def\@gobble#1{}
\long\def\@firstofone#1{#1}

% ---- trig (trig.sty) ---------------------------------------
\chardef\nin@ty=90
\chardef\@clxxx=180
\chardef\@lxxi=71
\mathchardef\@mmmmlxviii=4068
\chardef\@coeffz=72
\chardef\@coefb=42
\mathchardef\@coefc=840
\mathchardef\@coefd=5040
{\catcode`t=12\catcode`p=12\gdef\noPT#1pt{#1}}
\def\TG@rem@pt#1{\expandafter\noPT\the#1\space}
\def\TG@term#1{\dimen@\@tempb\dimen@\advance\dimen@ #1\p@}
\def\TG@series{%
  \dimen@\@lxxi\dimen@
  \divide\dimen@\@mmmmlxviii
  \edef\@tempa{\TG@rem@pt\dimen@}%
  \dimen@\@tempa\dimen@
  \edef\@tempb{\TG@rem@pt\dimen@}%
  \divide\dimen@\@coeffz
  \advance\dimen@\m@ne\p@
  \TG@term\@coefb
  \TG@term{-\@coefc}%
  \TG@term\@coefd
  \dimen@\@tempa\dimen@
  \divide\dimen@\@coefd}
\def\CalculateSin#1{{%
  \expandafter\ifx\csname sin(\number#1)\endcsname\relax
    \dimen@=#1\p@\TG@@sin
    \expandafter\xdef\csname sin(\number#1)\endcsname{\TG@rem@pt\dimen@}%
  \fi}}
\def\CalculateCos#1{{%
  \expandafter\ifx\csname cos(\number#1)\endcsname\relax
    \dimen@=\nin@ty\p@\advance\dimen@-#1\p@\TG@@sin
    \expandafter\xdef\csname cos(\number#1)\endcsname{\TG@rem@pt\dimen@}%
  \fi}}
\def\TG@reduce#1#2{%
  \dimen@#1#2\nin@ty\p@\advance\dimen@#2-\@clxxx\p@
  \dimen@-\dimen@\TG@@sin}
\def\TG@@sin{%
  \ifdim\TG@reduce>+%
  \else\ifdim\TG@reduce<-%
  \else\TG@series\fi\fi}
\def\UseSin#1{\csname sin(\number#1)\endcsname}
\def\UseCos#1{\csname cos(\number#1)\endcsname}
\def\z@num{0 }\def\@tempa{1 }\def\@tempb{-1 }
\expandafter\let\csname sin(0)\endcsname\z@num
\expandafter\let\csname cos(0)\endcsname\@tempa
\expandafter\let\csname sin(90)\endcsname\@tempa
\expandafter\let\csname cos(90)\endcsname\z@num
\expandafter\let\csname sin(-90)\endcsname\@tempb
\expandafter\let\csname cos(-90)\endcsname\z@num

% ---- number normalisation (pdftex.def) ---------------------
\def\GPT@space{ }
\def\GPT@MatrixIdentity{1 0 0 1}
\def\GPT@Zero{0}
\def\GPT@Minus{-}
\def\GPT@NormalizeNumber#1{%
  \edef#1{#1}%
  \edef#1{\expandafter\zap@space#1 \@empty}%
  \edef#1{\expandafter\GPT@ZapPlus#1+\@nil}%
  \edef#1{\expandafter\GPT@ZapMinusMinus#1--\@nil}%
  \expandafter\GPT@Split#1..\@nil
  \ifx\GPT@frac\@empty\else
    \edef\GPT@frac{\expandafter\GPT@Reverse\expandafter{\expandafter}\GPT@frac\@nil}%
    \edef\GPT@frac{\expandafter\GPT@ZapLeadingZeros\GPT@frac\@empty}%
    \ifx\GPT@frac\@empty\else
      \edef\GPT@frac{\expandafter\GPT@Reverse\expandafter{\expandafter}\GPT@frac\@nil}%
    \fi
  \fi
  \edef\GPT@sign{\expandafter\@car\GPT@int\@empty\@nil}%
  \ifx\GPT@sign\GPT@Minus
    \edef\GPT@int{\expandafter\@cdr\GPT@int\@nil}%
  \else\def\GPT@sign{}\fi
  \edef\GPT@int{\expandafter\GPT@ZapLeadingZeros\GPT@int\@empty}%
  \edef\GPT@temp{\GPT@int\GPT@frac}%
  \ifx\GPT@temp\@empty\def#1{0}%
  \else\edef#1{\GPT@sign\GPT@int\ifx\GPT@frac\@empty\else.\GPT@frac\fi}\fi}
\def\GPT@ZapPlus#1+#2\@nil{%
  #1\ifx\@empty#2\@empty\expandafter\@gobble\else\expandafter\@firstofone\fi
  {\GPT@ZapPlus#2\@nil}}
\def\GPT@ZapMinusMinus#1--#2\@nil{%
  #1\ifx\@empty#2\@empty\expandafter\@gobble\else\expandafter\@firstofone\fi
  {\GPT@ZapMinusMinus#2\@nil}}
\def\GPT@Split#1.#2.#3\@nil{%
  \def\GPT@int{#1}%
  \ifx\@empty#2\@empty\let\GPT@frac\@empty\else\def\GPT@frac{#2}\fi}
\def\GPT@Reverse#1#2#3\@nil{%
  \ifx\@empty#3\@empty#2#1\expandafter\@gobble
  \else\expandafter\@firstofone\fi
  {\GPT@Reverse{#2#1}#3\@nil}}
\def\GPT@ZapLeadingZeros#1{%
  \ifx0#1\expandafter\GPT@ZapLeadingZeros\else#1\fi}

% ---- bounding-box geometry (graphics.sty, unchanged) -------
\newdimen\Grot@height \newdimen\Grot@left  \newdimen\Grot@right
\newdimen\Grot@depth  \newdimen\Grot@l     \newdimen\Grot@r
\newdimen\Grot@h      \newdimen\Grot@d     \newdimen\Grot@x
\newdimen\Grot@y

\def\Grot@setangle#1{\edef\Grot@angle{#1}}
\def\Grot@Px#1#2#3{#1\Grot@cos#2\advance#1-\Grot@sin#3}
\def\Grot@Py#1#2#3{#1\Grot@sin#2\advance#1\Grot@cos#3}

\def\Grot@box{%
  \begingroup
  \CalculateSin\Grot@angle
  \CalculateCos\Grot@angle
  \edef\Grot@sin{\UseSin\Grot@angle}%
  \edef\Grot@cos{\UseCos\Grot@angle}%
  %
  % Measure the four corners of the original box relative to origin
  \Grot@r\wd\z@  \advance\Grot@r-\Grot@x   % right edge
  \Grot@l\z@     \advance\Grot@l-\Grot@x   % left edge
  \Grot@h\ht\z@  \advance\Grot@h-\Grot@y   % top edge
  \Grot@d-\dp\z@ \advance\Grot@d-\Grot@y   % bottom edge
  %
  % Compute rotated bounding box extents.
  % Which corners become the extremes depends on the quadrant of the angle.
  \ifdim\Grot@sin\p@>\z@
    \ifdim\Grot@cos\p@>\z@
      \Grot@Py\Grot@height\Grot@r\Grot@h   % top-right corner -> height
      \Grot@Px\Grot@right \Grot@r\Grot@d   % bottom-right     -> right
      \Grot@Px\Grot@left  \Grot@l\Grot@h   % top-left         -> left
      \Grot@Py\Grot@depth \Grot@l\Grot@d   % bottom-left      -> depth
    \else
      \Grot@Py\Grot@height\Grot@r\Grot@d
      \Grot@Px\Grot@right \Grot@l\Grot@d
      \Grot@Px\Grot@left  \Grot@r\Grot@h
      \Grot@Py\Grot@depth \Grot@l\Grot@h
    \fi
  \else
    \ifdim\Grot@cos\p@<\z@
      \Grot@Py\Grot@height\Grot@l\Grot@d
      \Grot@Px\Grot@right \Grot@l\Grot@h
      \Grot@Px\Grot@left  \Grot@r\Grot@d
      \Grot@Py\Grot@depth \Grot@r\Grot@h
    \else
      \Grot@Py\Grot@height\Grot@l\Grot@h
      \Grot@Px\Grot@right \Grot@r\Grot@h
      \Grot@Px\Grot@left  \Grot@l\Grot@d
      \Grot@Py\Grot@depth \Grot@r\Grot@d
    \fi
  \fi
  \advance\Grot@height\Grot@y
  \advance\Grot@depth\Grot@y
  %
  % Compute the shift needed to place the rotation origin correctly
  \Grot@Px\dimen@  \Grot@x\Grot@y
  \Grot@Py\dimen@ii\Grot@x\Grot@y
  \dimen@-\dimen@      \advance\dimen@-\Grot@left
  \dimen@ii-\dimen@ii  \advance\dimen@ii\Grot@y
  %
  % Build the inner box: kern to origin, raise to baseline,
  % apply PDF rotation transform, place content, close transform.
  % \Grot@start zeroes \wd\z@ so TeX does not advance by the
  % original (unrotated) width; the outer box below gets the
  % correct rotated width assigned explicitly.
  \setbox\z@\hbox{%
    \kern\dimen@
    \raise\dimen@ii\hbox{\Grot@start\box\z@\Grot@end}}%
  %
  % Stamp the rotated bounding box onto the outer wrapper box
  \ht\z@\Grot@height                          % rotated height above baseline
  \dp\z@-\Grot@depth                          % rotated depth below baseline
  \advance\Grot@right-\Grot@left              % total rotated width
  \wd\z@\Grot@right                           % <- this is what reserves space
  %                                           %    and prevents text collision
  \leavevmode\box\z@
  \endgroup}

% ---- LuaTeX PDF output -------------------------------------
%
%  \pdfextension literal writes raw bytes into the PDF content
%  stream without any current-point tracking, so:
%    - No "misplaced restore" warnings (unlike save/setmatrix/restore)
%    - \wd\z@\z@ is safe to use inside \Grot@start
%
%  Execution order inside \Grot@box's inner \setbox:
%    1. \Grot@start runs  -> emits  q [cos sin -sin cos 0 0] cm
%                         -> zeroes \wd of box 0 (the content box)
%    2. \box\z@   runs    -> places content with zero TeX width
%                            so TeX does not advance the pen
%    3. \Grot@end runs    -> emits  Q  (restores PDF graphics state)
%
%  The outer \setbox\z@ captures everything; \Grot@box then
%  overwrites its \ht, \dp, \wd with the computed rotated extents,
%  giving neighbouring text a proper box to flow around.
%
\def\Grot@start{%
  \GPT@NormalizeNumber\Grot@sin
  \GPT@NormalizeNumber\Grot@cos
  % Build the PDF rotation matrix: [cos sin -sin cos 0 0]
  \edef\GPT@temp{%
    \Grot@cos\GPT@space
    \Grot@sin\GPT@space
    \if-\Grot@sin\else
      \ifx\Grot@sin\GPT@Zero\GPT@Zero\else-\Grot@sin\fi
    \fi
    \GPT@space\Grot@cos}%
  \ifx\GPT@temp\GPT@MatrixIdentity
    % Angle is 0 degrees: no transform needed, no-op end
    \def\Grot@end{}%
  \else
    % Emit PDF save + rotation matrix
    \pdfextension literal {q \GPT@temp\GPT@space 0 0 cm}%
    % Zero the TeX-level width of the content box (box 0).
    % Without this, \box\z@ would advance the pen by the original
    % unrotated width, causing the next word to collide with the
    % rotated content. \Grot@box sets the correct rotated width
    % on the outer wrapper after this inner box is complete.
    \wd\z@\z@
  \fi}
\def\Grot@end{\pdfextension literal {Q}}

% ---- public interface --------------------------------------
\long\def\rotatebox#1#2{%
  \leavevmode
  \Grot@setangle{#1}%
  \setbox\z@\hbox{{#2}}%
  \Grot@x\z@
  \Grot@y\z@
  \Grot@box}
\catcode`@=12

\let\تدوير=\rotatebox
%%%%%%%%%%%%%%%%%%%%%%%%%%%%%%%%%%%%%%%%%%%%%%%%%%%%%%%%%%%%%%%%%%%%%%%%
%%%%%%%%%%%%%%%%%%%%%%%%%%%%%%%%%%%%%%%%%%%%%%%%%%%%%%%%%%%%%%%%%%%%%%%%
%%%%%%%%%%%%%%%%%%%%%%%%%%%%%%%%%%%%%%%%%%%%%%%%%%%%%%%%%%%%%%%%%%%%%%%%

\def\صورة#1#2#3#4{%
	\hbox{%
  \directlua{
  local converted1 = arabic_to_western_number(arabic_to_western_unit(reverse_digit_sequences('\luaescapestring{\unexpanded{#1}}')))
  local converted2 = arabic_to_western_number(arabic_to_western_unit(reverse_digit_sequences('\luaescapestring{\unexpanded{#2}}')))
  local converted3 = arabic_to_western_number(arabic_to_western_unit(reverse_digit_sequences('\luaescapestring{\unexpanded{#3}}')))
  local converted4 = reverse_digit_sequences('\luaescapestring{\unexpanded{#4}}')
  process_image(converted1,converted2,converted3, converted4)
  }% 
  }%
}

%%%%%%%%%%%%%%%%%

%%%%%%%%%%%%%%%%% comment:  صور %%%%%%%%%%%%%%%%%%%%%%%%%%%%%%%%%%%%%%%%%%%%%%
%%%%%%%%%%%%%%%%% comment:  صور %%%%%%%%%%%%%%%%%%%%%%%%%%%%%%%%%%%%%%%%%%%%%%

%%%%%%%%%%%%%%%%%%%%%%%%%%%%%%%%%%%%%%%%%%%%%%%%%%%%%%%%%%%%%%%%%%%%%%%%%%%%%%%%
%%%%%%%%%%%%%%%%% comment: Table of Contents %%%%%%%%%%%%%%%%%%%%%%%%%%%%%%%%%%
%%%%%%%%%%%%%%%%%%%%%%%%%%%%%%%%%%%%%%%%%%%%%%%%%%%%%%%%%%%%%%%%%%%%%%%%%%%%%%%%

%%%%%%%%%%%%%%%%% comment:  الفهرس %%%%%%%%%%%%%%%%%%%%%%%%%%%%%%%%%%%%%%%%%%%%%%
%%%%%%%%%%%%%%%%% comment:  الفهرس %%%%%%%%%%%%%%%%%%%%%%%%%%%%%%%%%%%%%%%%%%%%%%
\def\contentsname{\vskip 2em plus 1em minus 1em \سطروسط{\كبيرجدااا الفهرس} \vskip 1.5em plus 0.5em minus 0.5em}
\def\contents{%
\contentsname
  \directlua{
    local f = io.open("الفهرس.tex", "r")
    if f then
      f:close()
      tex.print("\string\\input{الفهرس.tex}")
    end
  }
  \directlua{ toc_open_for_overwrite() }%
}
%%%%%%%%%%%%%%%%%
\let\الفهرس=\contents

\def\سلام{%
  \directlua{toc_close()}
  \end
}
%%%%%%%%%%%%%%%%% comment:  الفهرس %%%%%%%%%%%%%%%%%%%%%%%%%%%%%%%%%%%%%%%%%%%%%%
%%%%%%%%%%%%%%%%% comment:  الفهرس %%%%%%%%%%%%%%%%%%%%%%%%%%%%%%%%%%%%%%%%%%%%%%

%%%%%%%%%%%%%%%%%%%%%%%%%%%%%%%%%%%%%%%%%%%%%%%%%%%%%%%%%%%%%%%%%%%%%%%%%%%%%%%%
%%%%%%%%%%%%%%%%% comment: Hyperlinks and Colors %%%%%%%%%%%%%%%%%%%%%%%%%%%%%%%
%%%%%%%%%%%%%%%%%%%%%%%%%%%%%%%%%%%%%%%%%%%%%%%%%%%%%%%%%%%%%%%%%%%%%%%%%%%%%%%%

%%%%%%%%%%%%%%%%% comment:  روابط %%%%%%%%%%%%%%%%%%%%%%%%%%%%%%%%%%%%%%%%%%%%%%
%%%%%%%%%%%%%%%%% comment:  روابط %%%%%%%%%%%%%%%%%%%%%%%%%%%%%%%%%%%%%%%%%%%%%%
% define an anchor at the current location
\def\hypertarget#1#2{%
  \pdfextension dest name {#1} xyz% declare named destination at this point
  #2%
}

% create a link to a named anchor
\def\hyperlink#1#2{%
  \leavevmode
  \changecolor{blue}{%
  \pdfextension startlink attr {/Border [0 0 0] /H /N} goto name {#1}%
    #2%
  \pdfextension endlink\relax}%
}

% external URL (open in browser)
\def\url#1#2{%
  \leavevmode
  \changecolor{blue}{%
  \pdfextension startlink attr {/Border [0 0 0] /H /N} user {/Subtype /Link /A << /S /URI /URI (#1) >>}%
    #2%
  \pdfextension endlink\relax}%
}

%%%%%%%%%%%%%%%%%
\let\مرجع=\hypertarget
\let\راجع=\hyperlink
\let\رابط=\url
%%%%%%%%%%%%%%%%% comment:  روابط %%%%%%%%%%%%%%%%%%%%%%%%%%%%%%%%%%%%%%%%%%%%%%
%%%%%%%%%%%%%%%%% comment:  روابط %%%%%%%%%%%%%%%%%%%%%%%%%%%%%%%%%%%%%%%%%%%%%%

%%%%%%%%%%%%%%%%%%%%%%%%%%%%%%%%%%%%%%%%%%%%%%%%%%%%%%%%%%%%%%%%%%%%%%%%%%%%%%%%
%%%%%%%%%%%%%%%%% comment: Colors %%%%%%%%%%%%%%%%%%%%%%%%%%%%%%%%%%%%%%%%%%%%%
%%%%%%%%%%%%%%%%%%%%%%%%%%%%%%%%%%%%%%%%%%%%%%%%%%%%%%%%%%%%%%%%%%%%%%%%%%%%%%%%

%%%%%%%%%%%%%%%%% comment:  ألوان %%%%%%%%%%%%%%%%%%%%%%%%%%%%%%%%%%%%%%%%%%%%%%
%%%%%%%%%%%%%%%%% comment:  ألوان %%%%%%%%%%%%%%%%%%%%%%%%%%%%%%%%%%%%%%%%%%%%%%
% Fixed version for LuaTeX compatibility
% Based on: https://tex.stackexchange.com/a/207307
% License: CC BY-SA 3.0

\newif\ifpdfmode

% Detect PDF mode (updated for LuaTeX compatibility)
\ifdefined\pdfoutput
  \ifnum\pdfoutput>0 %
    \pdfmodetrue
  \fi
\else
  \ifdefined\outputmode % LuaTeX specific
    \ifnum\outputmode=1 %
      \pdfmodetrue
    \fi
  \fi
\fi

% Define colors for PDF vs DVI mode
\edef\currentcolor{%
  \ifpdfmode
    0 g 0 G%
  \else
    gray 0%
  \fi
}



% Handle color stack for PDF vs DVI mode
\ifpdfmode
  % For PDF mode - use LuaTeX's \special{pdf:...} syntax
  \ifdefined\directlua
    % LuaTeX specific PDF commands
    \def\pdfcolorstackpush{\special{pdf: \currentcolor}}%
    \let\pdfcolorstackpop\pdfcolorstackpush
  \else
    % For pdfTeX
    \ifdefined\pdfcolorstack
      % Use pdfcolorstack if available
      \chardef\colorstackcnt=0 %
      \def\pdfcolorstackpush{%
        \pdfcolorstack\colorstackcnt push{\currentcolor}%
      }%
      \def\pdfcolorstackpop{%
        \pdfcolorstack\colorstackcnt pop\relax%
      }%
    \else
      % Fallback for older pdfTeX
      \def\pdfcolorstackpush{\pdfliteral{\currentcolor}}%
      \let\pdfcolorstackpop\pdfcolorstackpush
    \fi
  \fi
\else
  % For DVI mode only
  \def\pdfcolorstackpush{\special{color push \currentcolor}}%
  \def\pdfcolorstackpop{\special{color pop}}%
\fi

% Main color command
\edef\color#1{%
  \begingroup\noexpand\expandafter\noexpand\expandafter\noexpand\expandafter\endgroup
  \noexpand\expandafter\noexpand\ifx\noexpand\csname color#1\noexpand\endcsname\relax
    \def\noexpand\currentcolor{#1}%
  \noexpand\else
    \noexpand\expandafter\let\noexpand\expandafter\noexpand\currentcolor
      \noexpand\csname color#1\noexpand\endcsname
  \noexpand\fi
  \noexpand\pdfcolorstackpush
  \aftergroup\noexpand\resetcolor
}

% Reset color command
\edef\resetcolor{%
  \noexpand\pdfcolorstackpop
}

\edef\colorred{%
  \ifpdfmode
    1 0 0 rg 1 0 0 RG%
  \else
    rgb 1 0 0%
  \fi
}
\edef\colorblue{%
  \ifpdfmode
    0 0 1 rg 0 0 1 RG%
  \else
    rgb 0 0 1%
  \fi
}

\edef\colorgreen{%
  \ifpdfmode
    0 1 0 rg 0 1 0 RG%
  \else
    rgb 0 1 0%
  \fi
}
\def\changecolor#1#2{\hbox{\begingroup\color{#1}#2\endgroup}}


%%%%%%%%%%%%%%%%%
\let\colorأزرق=\colorblue
\let\colorأحمر=\colorred
\let\colorأخضر=\colorgreen
\let\لون=\changecolor
%%%%%%%%%%%%%%%%% comment:  ألوان %%%%%%%%%%%%%%%%%%%%%%%%%%%%%%%%%%%%%%%%%%%%%%
%%%%%%%%%%%%%%%%% comment:  ألوان %%%%%%%%%%%%%%%%%%%%%%%%%%%%%%%%%%%%%%%%%%%%%%

%%%%%%%%%%%%%%%%%%%%%%%%%%%%%%%%%%%%%%%%%%%%%%%%%%%%%%%%%%%%%%%%%%%%%%%%%%%%%%%%
%%%%%%%%%%%%%%%%% comment: Overset, Underset, Cancel, Custom Boxes %%%%%%%%%%%%%
%%%%%%%%%%%%%%%%%%%%%%%%%%%%%%%%%%%%%%%%%%%%%%%%%%%%%%%%%%%%%%%%%%%%%%%%%%%%%%%%
\catcode`\@=11

\newbox\binrelbox
\newbox\binrelrefbox
\let\binrel@@\mathord

\def\binrel@#1{%
  \setbox\binrelbox\hbox{$\m@th
    \thinmuskip0mu \medmuskip-1mu \thickmuskip0mu
    #1\mathrel{}$}%
  \setbox\binrelrefbox\hbox{$\m@th
    \thinmuskip0mu \medmuskip-1mu \thickmuskip0mu
    #1{}$}%
  \ifdim\wd\binrelrefbox<\wd\binrelbox
    \global\let\binrel@@\mathrel
  \else
    \setbox\binrelbox\hbox{$\m@th
      \thinmuskip0mu \medmuskip-1mu \thickmuskip0mu
      #1\mathbin{}$}%
    \ifdim\wd\binrelrefbox>\wd\binrelbox
      \global\let\binrel@@\mathbin
    \else
      \global\let\binrel@@\mathord
    \fi
  \fi
}

\def\overset#1#2{%
  \binrel@{#2}%
  \binrel@@{\mathop{\kern\z@ #2}\limits^{#1}}%
}

\def\underset#1#2{%
  \binrel@{#2}%
  \binrel@@{\mathop{\kern\z@ #2}\limits_{#1}}%
}

\catcode`\@=12
\let\فوق=\overset
\let\تحت=\underset
\def\اشطب#1{%
  \mathchoice
    {\cancelbox{\displaystyle #1\vphantom{ز}}}%
    {\cancelbox{\textstyle #1\vphantom{ز}}}%
    {\cancelbox{\scriptstyle #1\vphantom{ز}}}%
    {\cancelbox{\scriptscriptstyle #1 \vphantom{ز}}}%
}

\def\اشطب#1{%
  \mathchoice
    {\mathrel{\cancelbox{\displaystyle #1}}}%
    {\mathrel{\cancelbox{\textstyle #1}}}%
    {\mathrel{\cancelbox{\scriptstyle #1}}}%
    {\mathrel{\cancelbox{\scriptscriptstyle #1}}}%
}

\def\cancelbox#1{%
  \leavevmode
  \setbox0=\hbox{$#1$}%
  \copy0
  \pdfextension literal{\directlua{cancel()}}%
  \hskip-\wd0
  \hbox to \wd0{}%
}

%%%%%%%%%%%%%%%%%%%%%%%%%%%%%%%%%%%%%%%%%%%%%%%%%%%%%%%%%%%%%%%%%%%%%%%%%%%%%%%%
%%%%%%%%%%%%%%%%% comment: Custom Boxes with Optional Widths %%%%%%%%%%%%%%%%%%%
%%%%%%%%%%%%%%%%%%%%%%%%%%%%%%%%%%%%%%%%%%%%%%%%%%%%%%%%%%%%%%%%%%%%%%%%%%%%%%%%
\long\def\ahbox{\futurelet\ahboxnext\ahboxcheck}
\long\def\ahboxcheck{%
  \ifx\ahboxnext[%           % is the next token an open bracket?
    \let\ahboxnext\ahboxwith%   yes → switch to the width-aware handler
  \else%
    \let\ahboxnext\ahboxwithout% no  → switch to the plain handler
  \fi%
  \ahboxnext}                % call whichever handler was selected
\long\def\ahboxwith[#1]#2{%
\leavevmode%
\directlua{
local converted = arabic_to_western_number(arabic_to_western_unit(reverse_digit_sequences('\luaescapestring{\unexpanded{#1}}')))
local content = '\luaescapestring{\unexpanded{#2}}'
tex.sprint([=[\hbox to]=] .. converted .. [=[{]=] .. content.. [=[}]=])
}}
\long\def\ahboxwithout#1{\hbox{#1}}
\let\صندوق=\ahbox


%%%%%%%%%%%%%%%%%%%%%%%%%%%%%%%%%%%%%%%%%%%%%%%%%%%%%%%%%%%%%%%%%%%%%%%%%%%%%%%%%%%%%%%%
\long\def\avbox{\futurelet\avboxnext\avboxcheck}
\long\def\avboxcheck{%
  \ifx\avboxnext[%           % is the next token an open bracket?
    \let\avboxnext\avboxwith%   yes → switch to the width-aware handler
  \else%
    \let\avboxnext\avboxwithout% no  → switch to the plain handler
  \fi%
  \avboxnext}                % call whichever handler was selected
\long\def\avboxwith[#1]#2{%
\directlua{
local converted = arabic_to_western_number(arabic_to_western_unit(reverse_digit_sequences('\luaescapestring{\unexpanded{#1}}')))
local content = '\luaescapestring{\unexpanded{#2}}'
tex.sprint([=[\vbox to]=] .. converted .. [=[{]=] .. content.. [=[}]=])
}}
\long\def\avboxwithout#1{\vbox{#1}}
\let\صندوقة=\avbox
%%%%%%%%%%%%%%%%%%%%%%%%%%%%%%%%%%%%%%%%%%%%%%%%%%%%%%%%%%%%%%%%%%%%%%%%%%%%%%%%%%%%%%%%
\long\def\ahalign{\futurelet\ahalignnext\ahaligncheck}
\long\def\ahaligncheck{%
  \ifx\ahalignnext[%           % is the next token an open bracket?
    \let\ahalignnext\ahalignwith %   yes → switch to the width-aware handler
  \else
    \let\ahalignnext\ahalignwithout% no  → switch to the plain handler
  \fi
  \ahalignnext}                % call whichever handler was selected
\long\def\ahalignwith[#1]#2{
  \leavevmode
\directlua{
local converted = arabic_to_western_number(arabic_to_western_unit(reverse_digit_sequences('\luaescapestring{\unexpanded{#1}}')))
local content = '\luaescapestring{\unexpanded{#2}}'
tex.sprint([=[\halign to]=] .. converted .. [=[{]=] .. content.. [=[}]=])
}}
\long\def\ahalignwithout#1{\halign{#1}}
\let\حاذ=\ahalign
%%%%%%%%%%%%%%%%%%%%%%%%%%%%%%%%%%%%%%%%%%%%%%%%%%%%%%%%%%%%%%%%%%%%%%%%%%%%%%%%
%%%%%%%%%%%%%%%%% comment: Arabic Integrals (with reflection) %%%%%%%%%%%%%%%%%%
%%%%%%%%%%%%%%%%%%%%%%%%%%%%%%%%%%%%%%%%%%%%%%%%%%%%%%%%%%%%%%%%%%%%%%%%%%%%%%%%
\def\aint#1#2#3#4{\reflect{$\mathdir TLT \textdir TLT{#1∫_{\reflect{$\mathdir TRT \textdir TRT #2 {#3}$}}^{\reflect{$\mathdir TRT \textdir TRT #2 {#4}$}}}$}}
\def\aintegral#1#2{
	\mathop{
		\mathchoice
		{\aint{\displaystyle}{\scriptstyle}{#1}{#2}}
		{\aint{\textstyle}{\scriptstyle}{#1}{#2}}
		{\aint{\scriptstyle}{\scriptscriptstyle}{#1}{#2}}
		{\aint{\scriptscriptstyle}{\scriptscriptstyle}{#1}{#2}}
	}
}

% ------ Helper: test for optional argument ------
\def\arabicintegral{\futurelet\arabicintegralnext\arabicintegralcheck}
\def\arabicintegralcheck{%
  \ifx\arabicintegralnext[%
    \let\arabicintegralhandler\arabicintegralwith
  \else
    \let\arabicintegralhandler\arabicintegralwithout
  \fi
  \arabicintegralhandler
}
\def\arabicintegralwithout{\aintegral{}{}}

% ------ First optional argument exists ------
\def\arabicintegralwith[#1]{%
  \def\arabicintegralargone{#1}%
  \futurelet\arabicintegralnext\arabicintegralsecond
}
\def\arabicintegralsecond{%
  \ifx\arabicintegralnext[%
    \let\arabicintegralhandler\arabicintegralwithtwo
  \else
    \let\arabicintegralhandler\arabicintegralwithone
  \fi
  \arabicintegralhandler
}
\def\arabicintegralwithone{\aintegral{\arabicintegralargone}{}}
\def\arabicintegralwithtwo[#1]{\aintegral{\arabicintegralargone}{#1}}


\let\تكامل=\arabicintegral

%%%%%%%%%%%%%%%%%
\def\aoint#1#2#3#4{\reflect{$\mathdir TLT \textdir TLT{#1∮_{\reflect{$\mathdir TRT \textdir TRT #2 {#3}$}}^{\reflect{$\mathdir TRT \textdir TRT #2 {#4}$}}}$}}
\def\aointegral#1#2{
	\mathop{
		\mathchoice
		{\aoint{\displaystyle}{\scriptstyle}{#1}{#2}}
		{\aoint{\textstyle}{\scriptstyle}{#1}{#2}}
		{\aoint{\scriptstyle}{\scriptscriptstyle}{#1}{#2}}
		{\aoint{\scriptscriptstyle}{\scriptscriptstyle}{#1}{#2}}
	}
}
% ------ Helper: test for optional argument ------
\def\arabicointegral{\futurelet\arabicointegralnext\arabicointegralcheck}
\def\arabicointegralcheck{%
  \ifx\arabicointegralnext[%
    \let\arabicointegralhandler\arabicointegralwith
  \else
    \let\arabicointegralhandler\arabicointegralwithout
  \fi
  \arabicointegralhandler
}
\def\arabicointegralwithout{\aointegral{}{}}

% ------ First optional argument exists ------
\def\arabicointegralwith[#1]{%
  \def\arabicointegralargone{#1}%
  \futurelet\arabicointegralnext\arabicointegralsecond
}
\def\arabicointegralsecond{%
  \ifx\arabicointegralnext[%
    \let\arabicointegralhandler\arabicointegralwithtwo
  \else
    \let\arabicointegralhandler\arabicointegralwithone
  \fi
  \arabicointegralhandler
}
\def\arabicointegralwithone{\aointegral{\arabicointegralargone}{}}
\def\arabicointegralwithtwo[#1]{\aointegral{\arabicointegralargone}{#1}}


\let\تكاملة=\arabicointegral
%%%%%%%%%%%%%%%%%
\def\adoubleint#1#2#3#4{\reflect{$\mathdir TLT \textdir TLT{#1∬_{\reflect{$\mathdir TRT \textdir TRT #2 {#3}$}}^{\reflect{$\mathdir TRT \textdir TRT #2 {#4}$}}}$}}
\def\adoubleintegral#1#2{
	\mathop{
		\mathchoice
		{\adoubleint{\displaystyle}{\scriptstyle}{#1}{#2}}
		{\adoubleint{\textstyle}{\scriptstyle}{#1}{#2}}
		{\adoubleint{\scriptstyle}{\scriptscriptstyle}{#1}{#2}}
		{\adoubleint{\scriptscriptstyle}{\scriptscriptstyle}{#1}{#2}}
	}
}

% ------ Helper: test for optional argument ------
\def\arabicdoubleintegral{\futurelet\arabicdoubleintegralnext\arabicdoubleintegralcheck}
\def\arabicdoubleintegralcheck{%
  \ifx\arabicdoubleintegralnext[%
    \let\arabicdoubleintegralhandler\arabicdoubleintegralwith
  \else
    \let\arabicdoubleintegralhandler\arabicdoubleintegralwithout
  \fi
  \arabicdoubleintegralhandler
}
\def\arabicdoubleintegralwithout{\adoubleintegral{}{}}

% ------ First optional argument exists ------
\def\arabicdoubleintegralwith[#1]{%
  \def\arabicdoubleintegralargone{#1}%
  \futurelet\arabicdoubleintegralnext\arabicdoubleintegralsecond
}
\def\arabicdoubleintegralsecond{%
  \ifx\arabicdoubleintegralnext[%
    \let\arabicdoubleintegralhandler\arabicdoubleintegralwithtwo
  \else
    \let\arabicdoubleintegralhandler\arabicdoubleintegralwithone
  \fi
  \arabicdoubleintegralhandler
}
\def\arabicdoubleintegralwithone{\adoubleintegral{\arabicdoubleintegralargone}{}}
\def\arabicdoubleintegralwithtwo[#1]{\adoubleintegral{\arabicdoubleintegralargone}{#1}}


\let\تكاملان=\arabicdoubleintegral




%%%%%%%%%%%%%%%%%%%%%%%%%%%%%%%%%%%%%%%%%%%%%%%%%%%%%%%%%%%%%%%%%%%%%%%%%%%%%%%%
%%%%%%%%%%%%%%%%%%%%%%%%%%%%%%%%%%%%%%%%%%%%%%%%%%%%%%%%%%%%%%%%%%%%%%%%%%%%%%%%
%%%%%%%%%%%%%%%%%%%%%%%%%%%%%%%%%%%%%%%%%%%%%%%%%%%%%%%%%%%%%%%%%%%%%%%%%%%%%%%%
\def\alim#1#2{%
\mathop{%
  \setbox0=\hbox{$%
    \vcenter{\vbox{\hbox{$\scriptstyle #2$}\vskip1pt}}%
    \atop%
    \vcenter{\vbox{\vskip1pt\hbox{$\scriptstyle #1$}}}%
  $}%
  \setbox1=\hbox{ـ}%
  \setbox2=\hbox{ـا}%
  \setbox3=\hbox{نهـ}%
  \dimen0=\wd0%             width to cover
  \dimen1=\wd1%             width of one ـ
  \dimen2=0pt%              tatweel accumulator
  \count0=0%                tatweel counter
  \loop\ifdim\dimen2<\dimen0
    \advance\dimen2\dimen1
    \advance\count0 1%
  \repeat
  \dimen5=\wd3%
  \advance\dimen5\dimen2%
  \advance\dimen5\wd2%
  \advance\dimen2\wd2%
  \setbox4=\hbox{%
    نهـ%
    \count1=0%
    \loop\ifnum\count1<\count0
      ـ%
      \advance\count1 1%
    \repeat
    ـا%
    \kern-\dimen2%
    %\lower1pt\box0%
    \box0%
  }%
  \wd4=\dimen5%
  \box4%
}\nolimits}


\def\alimit#1#2{%
		\mathchoice
		{\alim{#1}{#2}}%
		{\alim{#1}{#2}}%
		{\mathop{\scalebox{0.7}{0.7}{$\alim{#1}{#2}$}}\nolimits}%
		{\mathop{\scalebox{0.5}{0.5}{$\alim{#1}{#2}$}}\nolimits}%
}


\def\arabiclimit#1{%
  \def\arabiclimitarg{#1}%        % save in a macro, not in input stream!
  \futurelet\arabiclimitnext\arabiclimitcheck%
}
\def\arabiclimitcheck{%
  \ifx\arabiclimitnext[%
    \expandafter\arabiclimitwith%
  \else
    \expandafter\arabiclimitwithout%
  \fi%
}
\def\arabiclimitwith[#1]{\alimit{\arabiclimitarg}{#1}}
\def\arabiclimitwithout{\alimit{\arabiclimitarg}{}}





\let\نهاية=\arabiclimit

%%%%%%%%%%%%%%%%%%%%%%%%%%%%%%%%%%%%%%%%%%%%%%%%%%%%%%%%%%%%%%%%%%%%%%%%%%%%%%%%
%%%%%%%%%%%%%%%%%%%%%%%%%%%%%%%%%%%%%%%%%%%%%%%%%%%%%%%%%%%%%%%%%%%%%%%%%%%%%%%%
%%%%%%%%%%%%%%%%%%%%%%%%%%%%%%%%%%%%%%%%%%%%%%%%%%%%%%%%%%%%%%%%%%%%%%%%%%%%%%%%
\def\asum#1#2{%
\mathop{%
  \setbox0=\hbox{$%
    \vcenter{\vbox{\hbox{$\scriptstyle #2$}\vskip1pt}}%
    \atop%
    \vcenter{\vbox{\vskip1pt\hbox{$\scriptstyle #1$}}}%
  $}%
  \setbox1=\hbox{ـ}%
  \setbox2=\hbox{ـ}%
  \setbox3=\hbox{مجـ}%
  \dimen0=\wd0%             width to cover
  \dimen1=\wd1%             width of one ـ
  \dimen2=0pt%              tatweel accumulator
  \count0=0%                tatweel counter
  \loop\ifdim\dimen2<\dimen0
    \advance\dimen2\dimen1
    \advance\count0 1%
  \repeat
  \dimen5=\wd3%
  \advance\dimen5\dimen2%
  \advance\dimen5\wd2%
  \advance\dimen2\wd2%
  \setbox4=\hbox{%
    مجـ%
    \count1=0%
    \loop\ifnum\count1<\count0
      ـ%
      \advance\count1 1%
    \repeat
    ـ%
    \kern-\dimen2%
    \lower1pt\box0%
  }%
  \wd4=\dimen5%
  \box4%
}\nolimits}

\def\asummation#1#2{%
		\mathchoice
		{\asum{#1}{#2}}%
		{\asum{#1}{#2}}%
		{\mathop{\scalebox{0.7}{0.7}{$\asum{#1}{#2}$}}\nolimits}%
		{\mathop{\scalebox{0.5}{0.5}{$\asum{#1}{#2}$}}\nolimits}%
}


\def\arabicsummation#1{%
  \def\arabicsummationarg{#1}%        % save in a macro, not in input stream!
  \futurelet\arabicsummationnext\arabicsummationcheck%
}
\def\arabicsummationcheck{%
  \ifx\arabicsummationnext[%
    \expandafter\arabicsummationwith%
  \else
    \expandafter\arabicsummationwithout%
  \fi%
}
\def\arabicsummationwith[#1]{\asummation{\arabicsummationarg}{#1}}
\def\arabicsummationwithout{\asummation{\arabicsummationarg}{}}


\let\مجموع=\arabicsummation



\def\amult#1#2{%
\mathop{%
  \setbox0=\hbox{$%
    \vcenter{\vbox{\hbox{$\scriptstyle #2$}\vskip1pt}}%
    \atop%
    \vcenter{\vbox{\vskip1pt\hbox{$\scriptstyle #1$}}}%
  $}%
  \setbox1=\hbox{ـ}%
  \setbox2=\hbox{ـ}%
  \setbox3=\hbox{مضـ}%
  \dimen0=\wd0%             width to cover
  \dimen1=\wd1%             width of one ـ
  \dimen2=0pt%              tatweel accumulator
  \count0=0%                tatweel counter
  \loop\ifdim\dimen2<\dimen0
    \advance\dimen2\dimen1
    \advance\count0 1%
  \repeat
  \dimen5=\wd3%
  \advance\dimen5\dimen2%
  \advance\dimen5\wd2%
  \advance\dimen2\wd2%
  \setbox4=\hbox{%
    مضـ%
    \count1=0%
    \loop\ifnum\count1<\count0
      ـ%
      \advance\count1 1%
    \repeat
    ـ%
    \kern-\dimen2%
    \lower1pt\box0%
  }%
  \wd4=\dimen5%
  \box4%
}\nolimits}

\def\amultiplication#1#2{%
		\mathchoice
		{\amult{#1}{#2}}%
		{\amult{#1}{#2}}%
		{\mathop{\scalebox{0.7}{0.7}{$\amult{#1}{#2}$}}\nolimits}%
		{\mathop{\scalebox{0.5}{0.5}{$\amult{#1}{#2}$}}\nolimits}%
}

\def\arabicmultiplication#1{%
  \def\arabicmultiplicationarg{#1}%        % save in a macro, not in input stream!
  \futurelet\arabicmultiplicationnext\arabicmultiplicationcheck%
}
\def\arabicmultiplicationcheck{%
  \ifx\arabicmultiplicationnext[%
    \expandafter\arabicmultiplicationwith%
  \else
    \expandafter\arabicmultiplicationwithout%
  \fi%
}
\def\arabicmultiplicationwith[#1]{\amultiplication{\arabicmultiplicationarg}{#1}}
\def\arabicmultiplicationwithout{\amultiplication{\arabicmultiplicationarg}{}}
\let\مضروب=\arabicmultiplication
%%%%%%%%%%%%%%%%%%%%%%%%%%%%%%%%%%%%%%%%%%%%%%%%%%%%%%%%%%%%%%%%%%%%%%%%%%%%%%%%
%%%%%%%%%%%%%%%%% comment: Input Buffer Callback (must be at end) %%%%%%%%%%%%%
%%%%%%%%%%%%%%%%%%%%%%%%%%%%%%%%%%%%%%%%%%%%%%%%%%%%%%%%%%%%%%%%%%%%%%%%%%%%%%%%

%%%%%%%%%%%%%%%%%%%%%%%%%%%%%%%%%%%%%%%%%%%%%%%%%%%%%%%%%%%%%%%%%%%%%%%%%%%%%%%%%%%%%%%%
%%%%%%%%%%%%%%%%%%%%%%%%%%%%%%%%%%%%%%%%%%%%%%%%%%%%%%%%%%%%%%%%%%%%%%%%%%%%%%%%%%%%%%%%
%%%%%%%%%%%%%%%%%%%%%%%%%%%%comment: مهم لا تغير مكان المعالج القبلي. يجب أن يكون في أخر الملف %%%%%%%%%%%%%%%%%%%%%%%%%%%%
%%%%%%%%%%%%%%%%%%%%%%%%%%%%%%%%%%%%%%%%%%%%%%%%%%%%%%%%%%%%%%%%%%%%%%%%%%%%%%%%%%%%%%%%
%%%%%%%%%%%%%%%%%%%%%%%%%%%%%%%%%%%%%%%%%%%%%%%%%%%%%%%%%%%%%%%%%%%%%%%%%%%%%%%%%%%%%%%%
\directlua{
  luatexbase.add_to_callback("process_input_buffer",
  function(line)
	return reverse_digit_sequences(line)
	end,
"change")
}
%%%%%%%%%%%%%%%%%%%%%%%%%%%%%%%%%%%%%%%%%%%%%%%%%%%%%%%%%%%%%%%%%%%%%%%%%%%%%%%%%%%%%%%%
%%%%%%%%%%%%%%%%%%%%%%%%%%%%%%%%%%%%%%%%%%%%%%%%%%%%%%%%%%%%%%%%%%%%%%%%%%%%%%%%%%%%%%%%
%%%%%%%%%%%%%%%%%%%%%%%%%%%%comment: مهم لا تغير مكان المعالج القبلي. يجب أن يكون في أخر الملف %%%%%%%%%%%%%%%%%%%%%%%%%%%%
%%%%%%%%%%%%%%%%%%%%%%%%%%%%%%%%%%%%%%%%%%%%%%%%%%%%%%%%%%%%%%%%%%%%%%%%%%%%%%%%%%%%%%%%
%%%%%%%%%%%%%%%%%%%%%%%%%%%%%%%%%%%%%%%%%%%%%%%%%%%%%%%%%%%%%%%%%%%%%%%%%%%%%%%%%%%%%%%%
